% !TeX program = lualatex
\documentclass[12pt,aspectratio=169]{beamer}
\usepackage[utf8]{inputenc}
\usepackage[T1]{fontenc}
\usepackage{amsmath,amsfonts,amssymb}
\usepackage{graphicx}
\usepackage{xcolor}
\usepackage{hyperref}
\usepackage{anyfontsize}
\usepackage{lmodern}
\usepackage{longtable}
\usepackage{array}
\usepackage{booktabs}

% ====== FONT SIZES ======
\renewcommand{\tiny}{\fontsize{4}{5}\selectfont}
\renewcommand{\scriptsize}{\fontsize{5}{6}\selectfont}
\renewcommand{\footnotesize}{\fontsize{6}{7}\selectfont}
\renewcommand{\small}{\fontsize{7}{8}\selectfont}
\renewcommand{\normalsize}{\fontsize{8}{9}\selectfont}
\renewcommand{\large}{\fontsize{9}{10}\selectfont}
\renewcommand{\Large}{\fontsize{10}{11}\selectfont}
\renewcommand{\LARGE}{\fontsize{11}{13}\selectfont}
\renewcommand{\huge}{\fontsize{12}{14}\selectfont}
\renewcommand{\Huge}{\fontsize{14}{17}\selectfont}

% ====== COLORS ======
\definecolor{redcorrect}{RGB}{200,30,30}
\definecolor{bluecorrect}{RGB}{30,80,200}
\definecolor{greencheck}{RGB}{0,150,50}
\definecolor{lightgray}{RGB}{245,245,245}
\definecolor{darkgray}{RGB}{50,50,50}
\definecolor{softyellow}{RGB}{255,248,200}
\definecolor{deepblue}{RGB}{20,60,150}
\definecolor{darkred}{RGB}{150,20,20}
\definecolor{orangewarning}{RGB}{255,140,0}
\definecolor{correctgreen}{RGB}{0,120,40}
\definecolor{trapcolor}{RGB}{180,80,0}

\setbeamercolor{title}{fg=white,bg=redcorrect}
\setbeamercolor{frametitle}{fg=white,bg=redcorrect}
\setbeamercolor{block title}{fg=white,bg=bluecorrect}
\setbeamercolor{block body}{fg=darkgray,bg=lightgray}
\setbeamercolor{normal text}{fg=darkgray}
\usecolortheme{default}

% ====== CUSTOM COMMANDS ======
\newcommand{\correct}[1]{\textcolor{greencheck}{\textbf{\checkmark}} #1}
\newcommand{\keypoint}[1]{\textcolor{deepblue}{\textbf{#1}}}
\newcommand{\answerbox}[1]{\fcolorbox{correctgreen}{green!8}{\parbox{0.88\textwidth}{\centering \textcolor{correctgreen}{\textbf{\checkmark\ CORRECT:}} #1}}}
\newcommand{\trapbox}[1]{\fcolorbox{trapcolor}{orange!8}{\parbox{0.88\textwidth}{\centering \textcolor{trapcolor}{\textbf{⚠ EXAM TRAP:}} #1}}}
\newcommand{\stepbox}[1]{%
	\begin{center}
		\fcolorbox{darkgray}{white}{%
			\parbox{0.90\textwidth}{\vspace{0.05cm} #1 \vspace{0.05cm}}%
		}
	\end{center}
}
\newcommand{\recall}[1]{\fcolorbox{deepblue}{blue!5}{\parbox{0.88\textwidth}{\centering \textbf{REMEMBER:} #1}}}

% ================================================================
% BEGIN DOCUMENT
% ================================================================
\begin{document}
	
	% ================================================================
	% TITLE SLIDE
	% ================================================================
	\begin{frame}
		\centering
		\vspace{1.2cm}
		{\fontsize{22}{26}\selectfont \textbf{Domain A: Understanding Financial Crime}}\\[0.3cm]
		{\fontsize{16}{19}\selectfont \textbf{Complete CAMS Study Guide}}\\[0.5cm]
		{\fontsize{11}{13}\selectfont \textit{All Questions from PNG Files — Fully Explained}}\\[0.3cm]
		\rule{5cm}{0.5pt}\\[0.2cm]
		{\fontsize{9}{11}\selectfont 
			\begin{tabular}{c}
				Money Laundering Stages \\
				Terrorist Financing vs Money Laundering \\
				Sanctions, PEPs, VASPs, DNFBPs \\
				Gatekeepers, Technology, Data \\
				FATF Recommendations \& Global Standards
			\end{tabular}
		}
		\vspace{0.5cm}
	\end{frame}
	
	% ================================================================
	% TABLE OF CONTENTS
	% ================================================================
	\begin{frame}
		\frametitle{\fontsize{13}{15}\selectfont \textbf{Table of Contents}}
		\vspace{0.2cm}
		\scriptsize
		\begin{longtable}{p{0.42\textwidth} p{0.53\textwidth}}
			\toprule
			\keypoint{\large Topic} & \keypoint{\large Questions} \\
			\midrule
			Money Laundering Definition \& Stages & Q1--Q2 \\
			Terrorist Financing vs Money Laundering & Q3--Q5 \\
			Sanctions: OFAC, Strict Liability, 50\% Rule & Q6--Q8 \\
			UK Bribery Act vs FCPA & Q9 \\
			Tax Evasion vs Tax Avoidance & Q10 \\
			Fraud Types: Synthetic Identity, Application & Q11--Q12 \\
			Predicate Offenses & Q13 \\
			Structuring / Smurfing & Q14--Q15 \\
			SAR vs STR — Who Files? & Q16--Q17 \\
			Tipping Off & Q18 \\
			Correspondent Banking & Q19--Q21 \\
			Trade-Based Money Laundering & Q22--Q23 \\
			Politically Exposed Persons (PEPs) & Q24--Q26 \\
			Beneficial Ownership & Q27 \\
			Customer Due Diligence (CDD) & Q28--Q29 \\
			MSB / PSP / E-Commerce & Q30--Q32 \\
			Insurance \& Casino ML Risks & Q33--Q34 \\
			Real Estate \& Art ML Risks & Q35--Q36 \\
			VASPs, Travel Rule, Privacy Coins & Q37--Q39 \\
			Securities Sector & Q40 \\
			Gatekeepers: Lawyers, TCSPs & Q41--Q42 \\
			Human Trafficking & Q43 \\
			Proliferation Financing & Q44 \\
			FATF Mutual Evaluations & Q45 \\
			Global AFC Standards \& Guidance & Q46--Q55 \\
			Technology for AFC Compliance & Q56--Q60 \\
			Data Collection \& Preparation & Q61--Q70 \\
			DNFBPs \& Free Trade Zones & Q71--Q75 \\
			ESG \& AML Integration & Q76--Q78 \\
			\bottomrule
		\end{longtable}
	\end{frame}
	
	% ================================================================
	% SECTION 1: ML DEFINITION & STAGES
	% ================================================================
	\begin{frame}
		\frametitle{\fontsize{13}{15}\selectfont \textbf{Section 1: Money Laundering Definition \& Stages}}
		\vspace{0.3cm}
		\begin{center}
			{\fontsize{12}{14}\selectfont \textbf{Understanding the Three Stages of ML}}
		\end{center}
		\vspace{0.2cm}
		\begin{block}{}
			\centering
			\large\textbf{Placement} $\rightarrow$ \textbf{Layering} $\rightarrow$ \textbf{Integration}
		\end{block}
		\vspace{0.1cm}
		\stepbox{\scriptsize
			\textbf{Placement:} Introducing illicit cash into the financial system (the most detectable stage). \\
			\textbf{Layering:} Obscuring origin through complex transactions (multiple accounts, jurisdictions). \\
			\textbf{Integration:} Returning ``cleaned'' funds to the legitimate economy (the ultimate goal).
		}
		\vspace{0.1cm}
		\recall{\scriptsize Money laundering is NOT just hiding money. The goal is \textbf{integration} so funds can be freely used.}
	\end{frame}
	
	% ================================================================
	% Q1: ML DEFINITION
	% ================================================================
	\begin{frame}
		\frametitle{\fontsize{12}{14}\selectfont \textbf{Q1: Money Laundering — Definition \& Scope}}
		\begin{block}{\fontsize{9}{11}\selectfont \textbf{Question}}
			\scriptsize
			A trainee states: \textit{``Money laundering just means hiding dirty money from authorities.''} \\
			Which response BEST corrects this understanding?
			\vspace{0.1cm}
			\begin{itemize}
				\item[\textbf{A.}] Correct — concealment from authorities is the primary objective.
				\item[\textbf{B.}] Incorrect — hiding is only the placement stage; the goal is integration so funds can be freely used.
				\item[\textbf{C.}] Partially correct — hiding is the goal but layering is what makes it criminal.
				\item[\textbf{D.}] Incorrect — ML is only illegal when the predicate offense is a felony.
			\end{itemize}
		\end{block}
		\vspace{0.05cm}
		\answerbox{\large \textbf{B} is correct.}
		\vspace{0.05cm}
		\stepbox{\scriptsize
			\keypoint{Explanation:} ML has three stages: \textbf{Placement}, \textbf{Layering}, and \textbf{Integration}. The ultimate goal is \textbf{INTEGRATION} — making funds usable without suspicion. FATF defines ML as \textit{``processing criminal proceeds to disguise their illegal origin.''} The trainee confuses hiding (placement only) with the full ML cycle.
		}
		\vspace{0.03cm}
		\trapbox{\scriptsize Hiding alone $\neq$ money laundering. Without \textbf{integration}, funds remain unusable.}
	\end{frame}
	
	% ================================================================
	% Q2: PLACEMENT vs LAYERING vs INTEGRATION
	% ================================================================
	\begin{frame}
		\frametitle{\fontsize{12}{14}\selectfont \textbf{Q2: Three Stages — Placement vs Layering vs Integration}}
		\begin{block}{\fontsize{9}{11}\selectfont \textbf{Question}}
			\scriptsize
			A criminal wires \$500K through 15 intermediary accounts in 6 countries before purchasing a real estate investment. Which ML stage BEST describes the wire transfer activity?
			\vspace{0.1cm}
			\begin{itemize}
				\item[\textbf{A.}] Placement — funds are entering international financial accounts.
				\item[\textbf{B.}] Layering — transactions designed to obscure the trail and origin of funds.
				\item[\textbf{C.}] Integration — funds are moving toward a legitimate asset purchase.
				\item[\textbf{D.}] Structuring — multiple transfers are used to stay below thresholds.
			\end{itemize}
		\end{block}
		\vspace{0.05cm}
		\answerbox{\large \textbf{B} is correct.}
		\vspace{0.05cm}
		\stepbox{\scriptsize
			\keypoint{Explanation:} \textbf{Layering} is characterized by complex multi-step transactions to create distance between illicit funds and their criminal origin. Multiple accounts, multiple jurisdictions, and rapid wire transfers are hallmark layering techniques. \textbf{Placement} = initial introduction of illicit cash. \textbf{Integration} = final stage where funds re-enter as apparently lawful assets.
		}
		\vspace{0.03cm}
		\trapbox{\scriptsize Multiple transfers $\neq$ structuring. Structuring = breaking up cash deposits to avoid CTR thresholds.}
	\end{frame}
	
	% ================================================================
	% SECTION 2: TF vs ML
	% ================================================================
	\begin{frame}
		\frametitle{\fontsize{13}{15}\selectfont \textbf{Section 2: Terrorist Financing vs Money Laundering}}
		\vspace{0.3cm}
		\begin{center}
			{\fontsize{12}{14}\selectfont \textbf{The Critical Distinction That CAMS Tests Most Frequently}}
		\end{center}
		\vspace{0.2cm}
		\begin{columns}[T]
			\begin{column}{0.48\textwidth}
				\begin{block}{\textbf{Money Laundering}}
					\scriptsize
					\begin{itemize}
						\item Funds are \textbf{ALWAYS ILLICIT}
						\item Focus: \textbf{SOURCE} of funds
						\item Goal: \textbf{HIDE ORIGIN}
						\item Three stages: Placement $\rightarrow$ Layering $\rightarrow$ Integration
					\end{itemize}
				\end{block}
			\end{column}
			\begin{column}{0.48\textwidth}
				\begin{block}{\textbf{Terrorist Financing}}
					\scriptsize
					\begin{itemize}
						\item Funds can be \textbf{100\% LEGITIMATE}
						\item Focus: \textbf{DESTINATION / INTENT}
						\item Goal: \textbf{FUND TERRORISM}
						\item No integration stage required — funds only need to reach the recipient
					\end{itemize}
				\end{block}
			\end{column}
		\end{columns}
		\vspace{0.1cm}
		\recall{\scriptsize \textbf{REMEMBER:} ML = where money came from (source). TF = where money is going (destination/purpose).}
	\end{frame}
	
	% ================================================================
	% Q3: TF vs ML
	% ================================================================
	\begin{frame}
		\frametitle{\fontsize{12}{14}\selectfont \textbf{Q3: Terrorist Financing vs Money Laundering — The Critical Distinction}}
		\begin{block}{\fontsize{9}{11}\selectfont \textbf{Question}}
			\scriptsize
			An investigator states: \textit{``This charity diverted legitimate donations to extremists — but since the money was legitimate, it cannot be terrorist financing.''} What is the critical error?
			\vspace{0.1cm}
			\begin{itemize}
				\item[\textbf{A.}] TF requires illegal source funds; legitimate funds cannot constitute TF.
				\item[\textbf{B.}] TF is defined by the PURPOSE of funds, not the source — legitimate funds diverted for terrorism IS TF.
				\item[\textbf{C.}] The error is that charity fraud, not TF, applies when source funds are legal.
				\item[\textbf{D.}] No error — TF requires both illegal source AND illegal destination.
			\end{itemize}
		\end{block}
		\vspace{0.05cm}
		\answerbox{\large \textbf{B} is correct.}
		\vspace{0.05cm}
		\stepbox{\scriptsize
			\keypoint{Explanation:} TF funds can be entirely \textbf{legitimate} (salaries, charity donations, savings) provided the \textbf{INTENT} is to support terrorism. FATF Recommendation 5 criminalizes TF regardless of source legality. The investigator confused source of funds (irrelevant for TF) with purpose/destination (determinative for TF).
		}
		\vspace{0.03cm}
		\trapbox{\scriptsize TF funds can be 100\% legitimate. \textbf{Destination and intent}, not source, are determinative.}
	\end{frame}
	
	% ================================================================
	% Q4: TF DETECTION
	% ================================================================
	\begin{frame}
		\frametitle{\fontsize{12}{14}\selectfont \textbf{Q4: Terrorist Financing — Why Hard to Detect}}
		\begin{block}{\fontsize{9}{11}\selectfont \textbf{Question}}
			\scriptsize
			Which characteristic makes terrorist financing transactions MOST DIFFICULT to detect compared to money laundering?
			\vspace{0.1cm}
			\begin{itemize}
				\item[\textbf{A.}] TF always uses cryptocurrency, which is untraceable.
				\item[\textbf{B.}] TF amounts are often small and below reporting thresholds, with no need to enter the legitimate financial system.
				\item[\textbf{C.}] TF is not criminalized in most jurisdictions, limiting detection cooperation.
				\item[\textbf{D.}] TF always involves foreign jurisdictions, making domestic detection impossible.
			\end{itemize}
		\end{block}
		\vspace{0.05cm}
		\answerbox{\large \textbf{B} is correct.}
		\vspace{0.05cm}
		\stepbox{\scriptsize
			\keypoint{Explanation:} TF transactions are often \textbf{small} (below CTR/STR reporting thresholds) and may involve entirely normal financial products. The funds do \textbf{NOT need to be ``cleaned''} — no integration stage is required. This makes TF behaviorally indistinguishable from normal transactions.
		}
		\vspace{0.03cm}
		\trapbox{\scriptsize TF $\neq$ large amounts. Small, legitimate-appearing transactions are the hallmark detection challenge.}
	\end{frame}
	
	% ================================================================
	% Q5: TF CRIMINALIZATION
	% ================================================================
	\begin{frame}
		\frametitle{\fontsize{12}{14}\selectfont \textbf{Q5: TF Criminalization — FATF Recommendation 5}}
		\begin{block}{\fontsize{9}{11}\selectfont \textbf{Question}}
			\scriptsize
			Under FATF Recommendation 5, which statement about the criminalization of terrorist financing is MOST ACCURATE?
			\vspace{0.1cm}
			\begin{itemize}
				\item[\textbf{A.}] TF must be criminalized only when the funds derive from drug trafficking.
				\item[\textbf{B.}] TF must be criminalized regardless of whether the funds are from legitimate or illicit sources, and regardless of whether a terrorist act actually occurs.
				\item[\textbf{C.}] TF criminalization is optional for countries that have ratified the UN Convention against Transnational Organized Crime.
				\item[\textbf{D.}] TF requires proof that funds were actually used to carry out a terrorist attack.
			\end{itemize}
		\end{block}
		\vspace{0.05cm}
		\answerbox{\large \textbf{B} is correct.}
		\vspace{0.05cm}
		\stepbox{\scriptsize
			\keypoint{Explanation:} FATF Recommendation 5 requires TF criminalization regardless of source and regardless of whether an attack actually occurs. The offense is complete when funds are provided/collected with terrorist intent.
		}
		\vspace{0.03cm}
		\trapbox{\scriptsize TF offense is complete on \textbf{intent and provision of funds}. No terrorist act needs to occur.}
	\end{frame}
	
	% ================================================================
	% SECTION 3: SANCTIONS
	% ================================================================
	\begin{frame}
		\frametitle{\fontsize{13}{15}\selectfont \textbf{Section 3: Sanctions — OFAC, Strict Liability, 50\% Rule}}
		\vspace{0.3cm}
		\begin{center}
			{\fontsize{12}{14}\selectfont \textbf{Key Sanctions Concepts for CAMS}}
		\end{center}
		\vspace{0.2cm}
		\begin{block}{\textbf{Three Critical Rules}}
			\scriptsize
			\begin{enumerate}
				\item \textbf{Strict Liability:} Civil sanctions violations require NO intent — ``We didn't know'' is NOT a defense.
				\item \textbf{50\% Rule:} Entities owned 50\%+ by a designated person are blocked even if not listed on the SDN list.
				\item \textbf{Comprehensive vs Targeted:} Goods screening $\neq$ sanctions screening — both must be checked independently.
			\end{enumerate}
		\end{block}
		\vspace{0.1cm}
		\recall{\scriptsize \textbf{REMEMBER:} SDN list screening alone is insufficient. Must trace ownership.}
	\end{frame}
	
	% ================================================================
	% Q6: OFAC 50% RULE
	% ================================================================
	\begin{frame}
		\frametitle{\fontsize{12}{14}\selectfont \textbf{Q6: Sanctions — OFAC 50\% Rule}}
		\begin{block}{\fontsize{9}{11}\selectfont \textbf{Question}}
			\scriptsize
			A compliance officer approves a payment to Company X because it is not on the OFAC SDN list. However, Company X is 60\% owned by a designated individual. What error has the officer made?
			\vspace{0.1cm}
			\begin{itemize}
				\item[\textbf{A.}] No error — only directly-listed entities are sanctioned under US law.
				\item[\textbf{B.}] The officer failed to apply OFAC's 50\% rule: entities 50\%+ owned by a designated person are blocked even if not themselves listed.
				\item[\textbf{C.}] The officer should have applied the EU's 25\% beneficial ownership threshold instead.
				\item[\textbf{D.}] The error is minor — a blocked transaction report satisfies compliance.
			\end{itemize}
		\end{block}
		\vspace{0.05cm}
		\answerbox{\large \textbf{B} is correct.}
		\vspace{0.05cm}
		\stepbox{\scriptsize
			\keypoint{Explanation:} OFAC's \textbf{50\% rule}: any entity \textbf{owned 50\% or more} by a designated person is itself blocked, regardless of whether it appears on the SDN list. Aggregate ownership counts (e.g., two designated persons each owning 30\% = 60\% combined = blocked).
		}
		\vspace{0.03cm}
		\trapbox{\scriptsize SDN list screening alone is \textbf{insufficient}. Must trace ownership.}
	\end{frame}
	
	% ================================================================
	% Q7: STRICT LIABILITY
	% ================================================================
	\begin{frame}
		\frametitle{\fontsize{12}{14}\selectfont \textbf{Q7: Sanctions — Strict Liability}}
		\begin{block}{\fontsize{9}{11}\selectfont \textbf{Question}}
			\scriptsize
			A bank processes a sanctioned transaction due to a data entry error, with no intent to evade sanctions. What is the correct legal position?
			\vspace{0.1cm}
			\begin{itemize}
				\item[\textbf{A.}] No violation — intent is required for a sanctions breach under OFAC regulations.
				\item[\textbf{B.}] Civil violation only — criminal penalties require intent, civil do not.
				\item[\textbf{C.}] Sanctions are strict liability — inadvertent violations can result in civil penalties even without intent.
				\item[\textbf{D.}] The bank is protected if it self-reports within 10 business days.
			\end{itemize}
		\end{block}
		\vspace{0.05cm}
		\answerbox{\large \textbf{C} is correct.}
		\vspace{0.05cm}
		\stepbox{\scriptsize
			\keypoint{Explanation:} Sanctions regimes are \textbf{strict liability for civil violations} — intent is NOT required. A bank can face civil penalties for inadvertent screening failures. Criminal prosecution typically requires intent/willfulness, but civil fines can be substantial even for honest mistakes.
		}
		\vspace{0.03cm}
		\trapbox{\scriptsize \textbf{Civil} sanctions violations = strict liability. ``We didn't know'' is NOT a defense.}
	\end{frame}
	
	% ================================================================
	% Q8: COMPREHENSIVE vs TARGETED
	% ================================================================
	\begin{frame}
		\frametitle{\fontsize{12}{14}\selectfont \textbf{Q8: Sanctions — Comprehensive vs Targeted}}
		\begin{block}{\fontsize{9}{11}\selectfont \textbf{Question}}
			\scriptsize
			A bank processes a payment involving goods not on any restricted list, between two non-sanctioned companies, but the beneficial owner of one company is a national of a comprehensively sanctioned country. The officer states: \textit{``The goods aren't restricted, so we can process this.''} Why is this wrong?
			\vspace{0.1cm}
			\begin{itemize}
				\item[\textbf{A.}] Correct reasoning — goods classification is the only relevant factor.
				\item[\textbf{B.}] Wrong — comprehensive country sanctions can prohibit all transactions involving nationals of that country.
				\item[\textbf{C.}] Wrong — the payment should be blocked only if over \$10,000.
				\item[\textbf{D.}] Correct only if the goods are dual-use items.
			\end{itemize}
		\end{block}
		\vspace{0.05cm}
		\answerbox{\large \textbf{B} is correct.}
		\vspace{0.05cm}
		\stepbox{\scriptsize
			\keypoint{Explanation:} \textbf{Comprehensive sanctions} (e.g., Iran, Cuba, North Korea) can block all transactions involving designated country nationals. \textbf{Goods/export controls} are a separate regime. The officer's error conflates goods controls with sanctions — both must be checked independently.
		}
		\vspace{0.03cm}
		\trapbox{\scriptsize Goods restriction lists $\neq$ sanctions lists. Sanctions screen \textbf{parties}; export controls screen \textbf{goods}.}
	\end{frame}
	
	% ================================================================
	% SECTION 4: UK BRIBERY ACT vs FCPA
	% ================================================================
	\begin{frame}
		\frametitle{\fontsize{13}{15}\selectfont \textbf{Section 4: UK Bribery Act vs FCPA}}
		\vspace{0.3cm}
		\begin{center}
			{\fontsize{12}{14}\selectfont \textbf{Facilitation Payments — The Critical Difference}}
		\end{center}
		\vspace{0.2cm}
		\begin{columns}[T]
			\begin{column}{0.48\textwidth}
				\begin{block}{\textbf{UK Bribery Act 2010}}
					\scriptsize
					\begin{itemize}
						\item \textbf{ZERO} facilitation payment tolerance
						\item \textbf{Global} extraterritorial reach
						\item Corporate offense: failure to prevent bribery
						\item Applies to UK companies worldwide
					\end{itemize}
				\end{block}
			\end{column}
			\begin{column}{0.48\textwidth}
				\begin{block}{\textbf{US FCPA}}
					\scriptsize
					\begin{itemize}
						\item Has a narrow facilitation payment exception
						\item Exception narrowed by DOJ/SEC in practice
						\item Many companies ban by policy
						\item Applies to US companies and issuers
					\end{itemize}
				\end{block}
			\end{column}
		\end{columns}
	\end{frame}
	
	% ================================================================
	% Q9: UK BRIBERY ACT vs FCPA
	% ================================================================
	\begin{frame}
		\frametitle{\fontsize{12}{14}\selectfont \textbf{Q9: UK Bribery Act vs FCPA — Facilitation Payments}}
		\begin{block}{\fontsize{9}{11}\selectfont \textbf{Question}}
			\scriptsize
			A UK-incorporated company operating in a country where local law permits facilitation payments makes such a payment. The regional manager argues local law makes it acceptable. The compliance officer disagrees. Who is correct?
			\vspace{0.1cm}
			\begin{itemize}
				\item[\textbf{A.}] Manager is correct — local law governs; UK Bribery Act defers to host country law.
				\item[\textbf{B.}] Both are wrong — all payments to officials are illegal under both UK and US law.
				\item[\textbf{C.}] Compliance officer is correct — UK Bribery Act 2010 prohibits facilitation payments globally, regardless of local law.
				\item[\textbf{D.}] Manager is correct — the FCPA's facilitation payment exception applies to UK companies.
			\end{itemize}
		\end{block}
		\vspace{0.05cm}
		\answerbox{\large \textbf{C} is correct.}
		\vspace{0.05cm}
		\stepbox{\scriptsize
			\keypoint{Explanation:} \textbf{UK Bribery Act 2010}: (1) Explicitly \textbf{prohibits} facilitation payments — NO exception; (2) \textbf{Extraterritorial reach} — applies to UK companies anywhere in the world. A UK company cannot use FCPA standards or local law to justify facilitation payments.
		}
		\vspace{0.03cm}
		\trapbox{\scriptsize UK Bribery Act = \textbf{ZERO} facilitation payment tolerance, \textbf{globally}.}
	\end{frame}
	
	% ================================================================
	% SECTION 5: TAX EVASION vs AVOIDANCE
	% ================================================================
	\begin{frame}
		\frametitle{\fontsize{13}{15}\selectfont \textbf{Section 5: Tax Evasion vs Tax Avoidance}}
		\vspace{0.3cm}
		\begin{center}
			{\fontsize{12}{14}\selectfont \textbf{The Critical Distinction Is Disclosure}}
		\end{center}
		\vspace{0.2cm}
		\begin{block}{}
			\centering
			\textbf{Tax Evasion} = Illegal concealment or misrepresentation of taxable income/assets. \\
			\vspace{0.1cm}
			\textbf{Tax Avoidance} = Legal use of legitimate structures to reduce tax burden. \\
			\vspace{0.1cm}
			\keypoint{The same offshore trust can be legal (avoidance) if fully disclosed, or illegal (evasion) if concealed.}
		\end{block}
	\end{frame}
	
	% ================================================================
	% Q10: TAX EVASION vs AVOIDANCE
	% ================================================================
	\begin{frame}
		\frametitle{\fontsize{12}{14}\selectfont \textbf{Q10: Tax Evasion vs Tax Avoidance}}
		\begin{block}{\fontsize{9}{11}\selectfont \textbf{Question}}
			\scriptsize
			A wealthy client uses offshore trusts and tax-haven structures. The key factor determining whether this is ILLEGAL tax evasion vs.\ LEGAL tax avoidance is:
			\vspace{0.1cm}
			\begin{itemize}
				\item[\textbf{A.}] Whether the jurisdiction has a tax treaty with the client's home country.
				\item[\textbf{B.}] Whether the structures were set up by a licensed attorney.
				\item[\textbf{C.}] Whether the client FULLY DISCLOSES all assets, income, and structures to relevant tax authorities.
				\item[\textbf{D.}] Whether the tax saved exceeds a materiality threshold.
			\end{itemize}
		\end{block}
		\vspace{0.05cm}
		\answerbox{\large \textbf{C} is correct.}
		\vspace{0.05cm}
		\stepbox{\scriptsize
			\keypoint{Explanation:} The \textbf{critical distinction is disclosure and truthfulness}. The same offshore trust can be legal (avoidance) if fully disclosed to authorities, or illegal (evasion) if concealed. Tax evasion is a \textbf{predicate offense for ML} in most FATF member jurisdictions.
		}
		\vspace{0.03cm}
		\trapbox{\scriptsize Offshore structures alone $\neq$ evasion. \textbf{Concealment from tax authorities} is the line.}
	\end{frame}
	
	% ================================================================
	% SECTION 6: FRAUD TYPES
	% ================================================================
	\begin{frame}
		\frametitle{\fontsize{13}{15}\selectfont \textbf{Section 6: Fraud Types — Synthetic Identity vs Identity Theft}}
		\vspace{0.3cm}
		\begin{center}
			{\fontsize{12}{14}\selectfont \textbf{Understanding the Difference}}
		\end{center}
		\vspace{0.2cm}
		\begin{columns}[T]
			\begin{column}{0.48\textwidth}
				\begin{block}{\textbf{Synthetic Identity Fraud}}
					\scriptsize
					\begin{itemize}
						\item \textbf{MANUFACTURED} identities
						\item Combines real elements (e.g., real SSN) with fabricated information
						\item No single victim reports the theft
						\item Harder to detect
						\item Bust-out pattern: build credit, extract maximum, disappear
					\end{itemize}
				\end{block}
			\end{column}
			\begin{column}{0.48\textwidth}
				\begin{block}{\textbf{Traditional Identity Theft}}
					\scriptsize
					\begin{itemize}
						\item \textbf{COMPLETE} stolen identity
						\item Uses a real, identified person's complete identity
						\item Victim reports the theft
						\item Detection often starts with victim complaint
					\end{itemize}
				\end{block}
			\end{column}
		\end{columns}
	\end{frame}
	
	% ================================================================
	% Q11: SYNTHETIC IDENTITY FRAUD
	% ================================================================
	\begin{frame}
		\frametitle{\fontsize{12}{14}\selectfont \textbf{Q11: Synthetic Identity Fraud vs Identity Theft}}
		\begin{block}{\fontsize{9}{11}\selectfont \textbf{Question}}
			\scriptsize
			A bank sees applications with different names and addresses but the SAME Social Security Number variants (slightly altered). This is BEST classified as:
			\vspace{0.1cm}
			\begin{itemize}
				\item[\textbf{A.}] Traditional identity theft — the SSN belongs to real victims being impersonated.
				\item[\textbf{B.}] Application fraud using stolen but unaltered identities.
				\item[\textbf{C.}] Synthetic identity fraud — new identities manufactured by combining real and fabricated data elements.
				\item[\textbf{D.}] Structuring — the applications are designed to stay below detection thresholds.
			\end{itemize}
		\end{block}
		\vspace{0.05cm}
		\answerbox{\large \textbf{C} is correct.}
		\vspace{0.05cm}
		\stepbox{\scriptsize
			\keypoint{Explanation:} \textbf{Synthetic identity fraud} creates MANUFACTURED identities by combining real elements (often a real SSN) with fabricated information (fake names, addresses, DOBs). Key detection signals: same SSN variant with multiple name/address combinations; no consistent credit history; rapid account opening followed by cash-out.
		}
		\vspace{0.03cm}
		\trapbox{\scriptsize Synthetic fraud $\neq$ identity theft. No single real victim — a new identity is \textbf{created} from mixed real/fake data.}
	\end{frame}
	
	% ================================================================
	% Q12: FRAUD PROCEEDS AS ML PREDICATE
	% ================================================================
	\begin{frame}
		\frametitle{\fontsize{12}{14}\selectfont \textbf{Q12: Fraud Proceeds as ML Predicate}}
		\begin{block}{\fontsize{9}{11}\selectfont \textbf{Question}}
			\scriptsize
			A bank identifies loan applications using fabricated employment and income documents. Loans are approved, quickly withdrawn, and transferred overseas. What TWO crimes have been committed?
			\vspace{0.1cm}
			\begin{itemize}
				\item[\textbf{A.}] Only fraud — money laundering requires a separate criminal organization.
				\item[\textbf{B.}] Only money laundering — the overseas transfer is the offense.
				\item[\textbf{C.}] Fraud (predicate offense) AND money laundering (laundering of fraud proceeds).
				\item[\textbf{D.}] Tax evasion and wire fraud only.
			\end{itemize}
		\end{block}
		\vspace{0.05cm}
		\answerbox{\large \textbf{C} is correct.}
		\vspace{0.05cm}
		\stepbox{\scriptsize
			\keypoint{Explanation:} Fraud is a \textbf{designated predicate offense} for money laundering under FATF standards. Both crimes are charged simultaneously: fraud (the underlying crime) and money laundering (laundering the proceeds). The bank must file a SAR.
		}
		\vspace{0.03cm}
		\trapbox{\scriptsize Fraud + overseas transfer = \textbf{two crimes}: fraud (predicate) AND ML (laundering proceeds).}
	\end{frame}
	
	% ================================================================
	% SECTION 7: PREDICATE OFFENSES
	% ================================================================
	\begin{frame}
		\frametitle{\fontsize{13}{15}\selectfont \textbf{Section 7: Predicate Offenses}}
		\vspace{0.3cm}
		\begin{center}
			{\fontsize{12}{14}\selectfont \textbf{FATF Designated Categories of Predicate Offenses}}
		\end{center}
		\vspace{0.2cm}
		\begin{block}{}
			\scriptsize
			FATF designates \textbf{21 categories of predicate offenses}, including:
			\vspace{0.1cm}
			\begin{columns}[T]
				\begin{column}{0.45\textwidth}
					\begin{itemize}
						\item Drug trafficking
						\item Corruption / bribery
						\item Fraud
						\item Counterfeiting
						\item Smuggling
						\item Human trafficking
					\end{itemize}
				\end{column}
				\begin{column}{0.45\textwidth}
					\begin{itemize}
						\item Tax crimes (in jurisdictions that criminalize them)
						\item Arms trafficking
						\item \textbf{Environmental crime}
						\item Cybercrime
						\item Terrorism
						\item Proliferation financing
					\end{itemize}
				\end{column}
			\end{columns}
			\vspace{0.1cm}
			\keypoint{Note:} Minor civil infractions (e.g., jaywalking, parking violations) are NOT predicate offenses.
		\end{block}
	\end{frame}
	
	% ================================================================
	% Q13: PREDICATE OFFENSES
	% ================================================================
	\begin{frame}
		\frametitle{\fontsize{12}{14}\selectfont \textbf{Q13: Predicate Offenses — FATF Categories}}
		\begin{block}{\fontsize{9}{11}\selectfont \textbf{Question}}
			\scriptsize
			Under FATF standards, which of the following is NOT a designated category of predicate offense for money laundering?
			\vspace{0.1cm}
			\begin{itemize}
				\item[\textbf{A.}] Drug trafficking.
				\item[\textbf{B.}] Tax crimes (in jurisdictions that criminalize them).
				\item[\textbf{C.}] Minor civil infractions (e.g., jaywalking, parking violations).
				\item[\textbf{D.}] Environmental crime.
			\end{itemize}
		\end{block}
		\vspace{0.05cm}
		\answerbox{\large \textbf{C} is correct.}
		\vspace{0.05cm}
		\stepbox{\scriptsize
			\keypoint{Explanation:} FATF designates \textbf{21 categories of predicate offenses} — all are \textit{serious crimes}. Minor civil infractions are NOT predicate offenses. Environmental crime has been an increasingly emphasized predicate in recent FATF guidance.
		}
		\vspace{0.03cm}
		\trapbox{\scriptsize Not ALL crimes are predicates — FATF focuses on \textbf{serious offenses}.}
	\end{frame}
	
	% ================================================================
	% SECTION 8: STRUCTURING / SMURFING
	% ================================================================
	\begin{frame}
		\frametitle{\fontsize{13}{15}\selectfont \textbf{Section 8: Structuring / Smurfing}}
		\vspace{0.3cm}
		\begin{center}
			{\fontsize{12}{14}\selectfont \textbf{The Crime That Does NOT Require Dirty Money}}
		\end{center}
		\vspace{0.2cm}
		\recall{\large \textbf{Structuring is illegal EVEN with clean money.}}
		\vspace{0.1cm}
		\begin{block}{}
			\centering
			\textbf{31 USC 5324:} The crime is deliberately breaking up cash transactions to keep them BELOW the reporting threshold (US: \$10,000 CTR threshold). \\
			\vspace{0.1cm}
			You do NOT need to prove the money is from crime to prove structuring. \\
			\vspace{0.1cm}
			Bank staff must file SARs and must NOT alert the customer (tipping off prohibition).
		\end{block}
	\end{frame}
	
	% ================================================================
	% Q14: STRUCTURING
	% ================================================================
	\begin{frame}
		\frametitle{\fontsize{12}{14}\selectfont \textbf{Q14: Structuring (Smurfing) — Key Rule}}
		\begin{block}{\fontsize{9}{11}\selectfont \textbf{Question}}
			\scriptsize
			A criminal uses 20 individuals to each deposit \$9,500 in cash at different branches of the same bank on the same day. Which statement is MOST accurate?
			\vspace{0.1cm}
			\begin{itemize}
				\item[\textbf{A.}] No crime — no single deposit exceeds \$10,000, so no CTR is required.
				\item[\textbf{B.}] Structuring (smurfing) — deliberately breaking up transactions to avoid the \$10,000 CTR threshold, which is a crime under 31 USC 5324 even if the underlying funds are entirely legitimate.
				\item[\textbf{C.}] Tax evasion — cash deposits are not ML if funds are from legal sources.
				\item[\textbf{D.}] Insider trading — multiple accounts used to conceal investment activity.
			\end{itemize}
		\end{block}
		\vspace{0.05cm}
		\answerbox{\large \textbf{B} is correct.}
		\vspace{0.05cm}
		\stepbox{\scriptsize
			\keypoint{Explanation:} \textbf{Structuring} is the illegal act of deliberately breaking up cash transactions to keep them BELOW a reporting threshold. \textbf{Critical}: Structuring is a crime under 31 USC 5324 \textbf{even if the underlying funds are entirely legitimate}. You do NOT need to prove the money is from crime. The criminal intent is avoiding the CTR filing.
		}
		\vspace{0.03cm}
		\trapbox{\scriptsize Structuring is illegal \textbf{even with clean money}. The crime is avoiding the reporting requirement.}
	\end{frame}
	
	% ================================================================
	% Q15: CTR vs SAR
	% ================================================================
	\begin{frame}
		\frametitle{\fontsize{12}{14}\selectfont \textbf{Q15: CTR vs SAR — When Each Is Required}}
		\begin{block}{\fontsize{9}{11}\selectfont \textbf{Question}}
			\scriptsize
			A US bank customer deposits exactly \$9,800 in cash every week for 6 months. No single deposit exceeds \$10,000. The bank has not filed CTRs but has filed a SAR. Which statement is MOST accurate?
			\vspace{0.1cm}
			\begin{itemize}
				\item[\textbf{A.}] Wrong — CTRs are required because the weekly pattern exceeds \$10,000.
				\item[\textbf{B.}] Partially correct — no CTR is required (no single transaction over \$10K), but the SAR is appropriate because the \$9,800 pattern indicates structuring.
				\item[\textbf{C.}] Wrong — the bank must file both CTRs and SARs for every \$9,800 deposit.
				\item[\textbf{D.}] Wrong — the bank's only obligation is to freeze the account.
			\end{itemize}
		\end{block}
		\vspace{0.05cm}
		\answerbox{\large \textbf{B} is correct.}
		\vspace{0.05cm}
		\stepbox{\scriptsize
			\keypoint{Explanation:} \textbf{CTRs} are required for SINGLE cash transactions over \$10,000. No individual deposit here exceeds \$10,000, so no CTR is legally required. However, the weekly \$9,800 pattern is textbook structuring — the bank is REQUIRED to file a \textbf{SAR}. CTRs are transaction-triggered; SARs are behavior/suspicion-triggered.
		}
		\vspace{0.03cm}
		\trapbox{\scriptsize \textbf{CTR} = triggered by single transactions over \$10K. \textbf{SAR} = triggered by suspicious behavior/patterns regardless of dollar amount.}
	\end{frame}
	
	% ================================================================
	% SECTION 9: SAR vs STR
	% ================================================================
	\begin{frame}
		\frametitle{\fontsize{13}{15}\selectfont \textbf{Section 9: SAR vs STR — Who Files?}}
		\vspace{0.3cm}
		\begin{center}
			{\fontsize{12}{14}\selectfont \textbf{The Most Frequently Tested Distinction}}
		\end{center}
		\vspace{0.2cm}
		\begin{columns}[T]
			\begin{column}{0.48\textwidth}
				\begin{block}{\textbf{US — SAR}}
					\scriptsize
					\begin{itemize}
						\item BSA obligation for \textbf{defined financial institutions}
						\item Banks, broker-dealers, MSBs, mutual funds, insurance companies, casinos
						\item US lawyers and real estate agents have \textbf{NO BSA SAR obligation}
						\item Filed with FinCEN
					\end{itemize}
				\end{block}
			\end{column}
			\begin{column}{0.48\textwidth}
				\begin{block}{\textbf{Non-US — STR}}
					\scriptsize
					\begin{itemize}
						\item Filed by DNFBPs and financial institutions
						\item Lawyers, real estate agents, casinos, TCSPs have STR obligations
						\item Filed with national FIU
						\item FATF Recommendation 23 covers DNFBPs
					\end{itemize}
				\end{block}
			\end{column}
		\end{columns}
		\vspace{0.1cm}
		\recall{\scriptsize \textbf{REMEMBER:} US lawyers/real estate agents do NOT file SARs. Non-US DNFBPs file STRs with national FIUs.}
	\end{frame}
	
	% ================================================================
	% Q16: SAR vs STR
	% ================================================================
	\begin{frame}
		\frametitle{\fontsize{12}{14}\selectfont \textbf{Q16: SAR vs STR — Who Has the Obligation?}}
		\begin{block}{\fontsize{9}{11}\selectfont \textbf{Question}}
			\scriptsize
			Under the US Bank Secrecy Act (BSA), which entities are REQUIRED to file Suspicious Activity Reports (SARs)?
			\vspace{0.1cm}
			\begin{itemize}
				\item[\textbf{A.}] All US businesses whenever they suspect a crime involving money.
				\item[\textbf{B.}] Only federally-chartered banks.
				\item[\textbf{C.}] BSA-defined ``financial institutions'': banks, broker-dealers, MSBs, mutual funds, insurance companies, casinos — NOT US attorneys or US real estate agents.
				\item[\textbf{D.}] All FATF-member entities universally.
			\end{itemize}
		\end{block}
		\vspace{0.05cm}
		\answerbox{\large \textbf{C} is correct.}
		\vspace{0.05cm}
		\stepbox{\scriptsize
			\keypoint{Explanation:} In the US, SAR filing is a BSA obligation for \textbf{defined financial institutions} only. US \textbf{attorneys} and \textbf{real estate agents} do NOT have BSA SAR obligations. In other jurisdictions, DNFBPs file STRs with national FIUs.
		}
		\vspace{0.03cm}
		\trapbox{\scriptsize US SAR = BSA obligation for \textbf{defined FIs}. US lawyers/real estate agents $\neq$ BSA SAR filers.}
	\end{frame}
	
	% ================================================================
	% Q17: STR OBLIGATION — DNFBPs
	% ================================================================
	\begin{frame}
		\frametitle{\fontsize{12}{14}\selectfont \textbf{Q17: STR Obligation — DNFBPs Outside the US}}
		\begin{block}{\fontsize{9}{11}\selectfont \textbf{Question}}
			\scriptsize
			In a FATF-member jurisdiction (outside the US), a lawyer assists with a real estate transaction and identifies suspicious circumstances. What AML obligation does the lawyer typically have?
			\vspace{0.1cm}
			\begin{itemize}
				\item[\textbf{A.}] No obligation — attorney-client privilege fully exempts lawyers from AML reporting.
				\item[\textbf{B.}] File a SAR with FinCEN if the client is American.
				\item[\textbf{C.}] File a Suspicious Transaction Report (STR) with the national FIU for designated activities, subject to legal professional privilege carve-outs.
				\item[\textbf{D.}] Report directly to the client's home country regulators.
			\end{itemize}
		\end{block}
		\vspace{0.05cm}
		\answerbox{\large \textbf{C} is correct.}
		\vspace{0.05cm}
		\stepbox{\scriptsize
			\keypoint{Explanation:} Under FATF Recommendation 23, DNFBPs including lawyers are required to file STRs when they perform \textbf{designated activities}: real estate transactions, managing client funds, company/trust formation. \textbf{Legal professional privilege} protects legal \textit{advice} but does NOT protect transactional facilitation.
		}
		\vspace{0.03cm}
		\trapbox{\scriptsize Lawyers DO have STR obligations for \textbf{designated transactional activities}. LPP does NOT protect transactional facilitation.}
	\end{frame}
	
	% ================================================================
	% SECTION 10: TIPPING OFF
	% ================================================================
	\begin{frame}
		\frametitle{\fontsize{13}{15}\selectfont \textbf{Section 10: Tipping Off}}
		\vspace{0.3cm}
		\begin{center}
			{\fontsize{12}{14}\selectfont \textbf{A Criminal Offense in All FATF Jurisdictions}}
		\end{center}
		\vspace{0.2cm}
		\recall{\large \textbf{Tipping off is a CRIMINAL OFFENSE.}}
		\vspace{0.1cm}
		\begin{block}{}
			\scriptsize
			Once a SAR/STR has been filed, staff are PROHIBITED from disclosing to the subject:
			\begin{itemize}
				\item That a SAR/STR has been filed
				\item That an ML/TF investigation is underway
			\end{itemize}
			\vspace{0.1cm}
			\textbf{US:} 31 USC 5318(g)(2) — SAR secrecy is absolute. \\
			\textbf{Exception:} Intra-group sharing in certain jurisdictions (limited).
		\end{block}
	\end{frame}
	
	% ================================================================
	% Q18: TIPPING OFF
	% ================================================================
	\begin{frame}
		\frametitle{\fontsize{12}{14}\selectfont \textbf{Q18: Tipping Off — Criminal Prohibition}}
		\begin{block}{\fontsize{9}{11}\selectfont \textbf{Question}}
			\scriptsize
			A bank files a SAR on a customer suspected of ML. The customer's relationship manager then warns the customer that a report has been filed. What has occurred?
			\vspace{0.1cm}
			\begin{itemize}
				\item[\textbf{A.}] A minor procedural breach — the SAR remains valid but must be refiled.
				\item[\textbf{B.}] The tipping-off offense — illegal disclosure of a SAR filing to the subject, carrying criminal penalties.
				\item[\textbf{C.}] A legitimate customer service action — customers have a right to know when they are reported.
				\item[\textbf{D.}] A GDPR/data privacy breach only — no AML-specific offense.
			\end{itemize}
		\end{block}
		\vspace{0.05cm}
		\answerbox{\large \textbf{B} is correct.}
		\vspace{0.05cm}
		\stepbox{\scriptsize
			\keypoint{Explanation:} \textbf{Tipping off} is a criminal offense. Once a SAR/STR has been filed, staff are PROHIBITED from disclosing to the subject that a SAR has been filed or that an investigation is underway. Staff face criminal prosecution.
		}
		\vspace{0.03cm}
		\trapbox{\scriptsize Tipping off = \textbf{criminal offense}. SAR secrecy is absolute toward the subject.}
	\end{frame}
	
	% ================================================================
	% SECTION 11: CORRESPONDENT BANKING
	% ================================================================
	\begin{frame}
		\frametitle{\fontsize{13}{15}\selectfont \textbf{Section 11: Correspondent Banking}}
		\vspace{0.3cm}
		\begin{center}
			{\fontsize{12}{14}\selectfont \textbf{FATF Recommendation 13 — Key Requirements}}
		\end{center}
		\vspace{0.2cm}
		\begin{block}{}
			\scriptsize
			\textbf{EDD Required:}
			\begin{itemize}
				\item Gather information on respondent's business, reputation, and AML supervision
				\item Assess respondent's AML/CFT controls
				\item Obtain \textbf{senior management approval}
				\item Document respective AML responsibilities
				\item Ensure respondent does \textbf{not permit shell banks}
			\end{itemize}
			\vspace{0.1cm}
			\keypoint{FATF explicitly discourages blanket de-risking.}
		\end{block}
	\end{frame}
	
	% ================================================================
	% Q19: CORRESPONDENT BANKING
	% ================================================================
	\begin{frame}
		\frametitle{\fontsize{12}{14}\selectfont \textbf{Q19: Correspondent Banking — FATF Recommendation 13}}
		\begin{block}{\fontsize{9}{11}\selectfont \textbf{Question}}
			\scriptsize
			A US bank reviews a correspondent relationship with a respondent bank in a high-risk jurisdiction with weak AML controls. Under FATF Recommendation 13, what is the US bank's PRIMARY obligation?
			\vspace{0.1cm}
			\begin{itemize}
				\item[\textbf{A.}] Terminate all correspondent relationships with high-risk jurisdictions immediately.
				\item[\textbf{B.}] Conduct EDD on the respondent: assess AML controls, understand ownership, verify no shell bank clients, and obtain senior management approval.
				\item[\textbf{C.}] Limit the relationship to USD correspondent services only.
				\item[\textbf{D.}] File SARs covering all transactions for the prior 12 months.
			\end{itemize}
		\end{block}
		\vspace{0.05cm}
		\answerbox{\large \textbf{B} is correct.}
		\vspace{0.05cm}
		\stepbox{\scriptsize
			\keypoint{Explanation:} FATF Recommendation 13 requires EDD including senior management approval. Outright termination (``de-risking'') is \textbf{NOT FATF's preferred approach} — FATF explicitly discourages blanket de-risking. The obligation is to \textit{manage} risk, not reflexively exit.
		}
		\vspace{0.03cm}
		\trapbox{\scriptsize FATF \textbf{opposes blanket de-risking}. Senior management approval is mandatory.}
	\end{frame}
	
	% ================================================================
	% Q20: SHELL BANKS
	% ================================================================
	\begin{frame}
		\frametitle{\fontsize{12}{14}\selectfont \textbf{Q20: Shell Banks — Absolute Prohibition}}
		\begin{block}{\fontsize{9}{11}\selectfont \textbf{Question}}
			\scriptsize
			FATF Recommendation 13 explicitly prohibits banks from doing what regarding shell banks?
			\vspace{0.1cm}
			\begin{itemize}
				\item[\textbf{A.}] Offering currency exchange services to shell banks.
				\item[\textbf{B.}] Establishing or maintaining correspondent relationships with shell banks, or allowing respondent banks to use their accounts to serve shell banks.
				\item[\textbf{C.}] Accepting deposits from companies registered in tax haven jurisdictions.
				\item[\textbf{D.}] Conducting wire transfers for any entity incorporated offshore.
			\end{itemize}
		\end{block}
		\vspace{0.05cm}
		\answerbox{\large \textbf{B} is correct.}
		\vspace{0.05cm}
		\stepbox{\scriptsize
			\keypoint{Explanation:} FATF Recommendation 13 explicitly prohibits correspondent relationships with \textbf{shell banks}. A \textbf{shell bank} = a bank with \textit{no physical presence} in any country and \textit{not affiliated} with a regulated financial group. This prohibition is \textbf{ABSOLUTE} — no amount of EDD makes it acceptable.
		}
		\vspace{0.03cm}
		\trapbox{\scriptsize Shell bank prohibition is \textbf{ABSOLUTE} — no risk-based exception exists.}
	\end{frame}
	
	% ================================================================
	% Q21: NESTED CORRESPONDENT ACCOUNTS
	% ================================================================
	\begin{frame}
		\frametitle{\fontsize{12}{14}\selectfont \textbf{Q21: Nested Correspondent Accounts}}
		\begin{block}{\fontsize{9}{11}\selectfont \textbf{Question}}
			\scriptsize
			What is a ``nested correspondent account'' and why is it particularly high risk?
			\vspace{0.1cm}
			\begin{itemize}
				\item[\textbf{A.}] A correspondent account used only for domestic wire transfers.
				\item[\textbf{B.}] An account where a third-party institution accesses the correspondent bank's services INDIRECTLY through the respondent, without the correspondent knowing the ultimate end-users.
				\item[\textbf{C.}] A savings account nested within a checking account for cash management.
				\item[\textbf{D.}] A correspondent relationship between two banks in the same jurisdiction.
			\end{itemize}
		\end{block}
		\vspace{0.05cm}
		\answerbox{\large \textbf{B} is correct.}
		\vspace{0.05cm}
		\stepbox{\scriptsize
			\keypoint{Explanation:} A \textbf{nested correspondent account} occurs when a respondent bank allows ITS OWN CLIENTS to conduct transactions through the respondent's correspondent account. The correspondent has done CDD on the respondent but has NO VISIBILITY into the respondent's clients. This creates a ``black box'' for illicit actors.
		}
		\vspace{0.03cm}
		\trapbox{\scriptsize Nested accounts = correspondent's services reach \textbf{unknown end-users}.}
	\end{frame}
	
	% ================================================================
	% SECTION 12: TBML
	% ================================================================
	\begin{frame}
		\frametitle{\fontsize{13}{15}\selectfont \textbf{Section 12: Trade-Based Money Laundering}}
		\vspace{0.3cm}
		\begin{center}
			{\fontsize{12}{14}\selectfont \textbf{Using Trade to Transfer Illicit Value}}
		\end{center}
		\vspace{0.2cm}
		\begin{block}{}
			\centering
			\textbf{Over-invoicing:} Buyer pays more than goods are worth $\rightarrow$ value moves TO seller. \\
			\textbf{Under-invoicing:} Buyer pays less than goods are worth $\rightarrow$ value moves TO buyer. \\
			\vspace{0.1cm}
			Value transferred = (billed price $-$ market price) $\times$ quantity.
		\end{block}
		\vspace{0.1cm}
		\recall{\scriptsize TBML is extremely difficult to detect without trade finance expertise and commodity market knowledge.}
	\end{frame}
	
	% ================================================================
	% Q22: TBML
	% ================================================================
	\begin{frame}
		\frametitle{\fontsize{12}{14}\selectfont \textbf{Q22: Trade-Based Money Laundering (TBML)}}
		\begin{block}{\fontsize{9}{11}\selectfont \textbf{Question}}
			\scriptsize
			An invoice shows 10,000 tons of steel at \$500/ton when the market price is \$250/ton. What ML method does this represent, and what is the primary purpose?
			\vspace{0.1cm}
			\begin{itemize}
				\item[\textbf{A.}] Under-invoicing — to avoid import duties on the steel.
				\item[\textbf{B.}] Over-invoicing — to transfer \$2.5M in illicit value from the buyer's jurisdiction to the seller's jurisdiction.
				\item[\textbf{C.}] Phantom shipment — no goods are actually shipped.
				\item[\textbf{D.}] Multiple invoicing — the same shipment is invoiced to different buyers.
			\end{itemize}
		\end{block}
		\vspace{0.05cm}
		\answerbox{\large \textbf{B} is correct.}
		\vspace{0.05cm}
		\stepbox{\scriptsize
			\keypoint{Explanation:} \textbf{Over-invoicing}: buyer pays more than goods are worth. Calculation: \$500 billed vs \$250 market $\rightarrow$ \$250/ton $\times$ 10,000 tons = \textbf{\$2.5M illicit value transferred to the seller}. \textbf{Under-invoicing} moves value to the buyer.
		}
		\vspace{0.03cm}
		\trapbox{\scriptsize \textbf{Over-invoicing} moves value TO the seller. \textbf{Under-invoicing} moves value TO the buyer.}
	\end{frame}
	
	% ================================================================
	% Q23: FREE TRADE ZONES
	% ================================================================
	\begin{frame}
		\frametitle{\fontsize{12}{14}\selectfont \textbf{Q23: Free Trade Zones (FTZ) — ML Risk}}
		\begin{block}{\fontsize{9}{11}\selectfont \textbf{Question}}
			\scriptsize
			Why are Free Trade Zones specifically identified by FATF as HIGH RISK for trade-based money laundering?
			\vspace{0.1cm}
			\begin{itemize}
				\item[\textbf{A.}] FTZs operate outside national legal systems making contracts unenforceable.
				\item[\textbf{B.}] FTZs only accept payments in foreign currency, complicating AML screening.
				\item[\textbf{C.}] FTZs often have reduced customs inspection, minimal AML oversight, and relaxed documentation requirements.
				\item[\textbf{D.}] FTZs require all transactions to use a single currency.
			\end{itemize}
		\end{block}
		\vspace{0.05cm}
		\answerbox{\large \textbf{C} is correct.}
		\vspace{0.05cm}
		\stepbox{\scriptsize
			\keypoint{Explanation:} FATF identified FTZs as high-risk due to: (1) \textbf{Reduced inspection}; (2) \textbf{Documentation gaps}; (3) \textbf{Minimal AML oversight}; (4) \textbf{Multiple company layering}. Classic TBML in FTZ: invoice a \$10 widget as \$1,000, ship through FTZ with minimal inspection.
		}
		\vspace{0.03cm}
		\trapbox{\scriptsize FTZ risk = \textbf{reduced inspection + minimal AML oversight + documentation gaps}.}
	\end{frame}
	
	% ================================================================
	% SECTION 13: PEPs
	% ================================================================
	\begin{frame}
		\frametitle{\fontsize{13}{15}\selectfont \textbf{Section 13: Politically Exposed Persons (PEPs)}}
		\vspace{0.3cm}
		\begin{center}
			{\fontsize{12}{14}\selectfont \textbf{FATF Recommendation 12 — Key Requirements}}
		\end{center}
		\vspace{0.2cm}
		\begin{block}{}
			\scriptsize
			\textbf{PEP Definition:} Individuals entrusted with prominent public functions (heads of state, senior politicians, senior government/judicial/military officials) \textbf{AND their family members and close associates}.
			\vspace{0.1cm}
			\textbf{EDD Requirements:}
			\begin{itemize}
				\item \textbf{Senior management approval} to establish/continue the relationship
				\item \textbf{Source of Wealth} AND \textbf{Source of Funds} — BOTH required
				\item \textbf{Enhanced ongoing monitoring}
			\end{itemize}
		\end{block}
	\end{frame}
	
	% ================================================================
	% Q24: PEP DEFINITION
	% ================================================================
	\begin{frame}
		\frametitle{\fontsize{12}{14}\selectfont \textbf{Q24: PEP Definition — FATF Recommendation 12}}
		\begin{block}{\fontsize{9}{11}\selectfont \textbf{Question}}
			\scriptsize
			Under FATF Recommendation 12, which individual is classified as a PEP?
			\vspace{0.1cm}
			\begin{itemize}
				\item[\textbf{A.}] A retired local mayor who left office 10 years ago with no remaining influence.
				\item[\textbf{B.}] A mid-level customs officer with no significant discretionary authority.
				\item[\textbf{C.}] The adult child of a current foreign head of state who runs a private business.
				\item[\textbf{D.}] A senior bank compliance officer at a regulated financial institution.
			\end{itemize}
		\end{block}
		\vspace{0.05cm}
		\answerbox{\large \textbf{C} is correct.}
		\vspace{0.05cm}
		\stepbox{\scriptsize
			\keypoint{Explanation:} FATF Recommendation 12 defines PEPs as individuals entrusted with \textbf{prominent public functions} \textbf{AND their family members and close associates}. The adult child of a current foreign head of state = PEP (family member category).
		}
		\vspace{0.03cm}
		\trapbox{\scriptsize PEP family members and close associates are \textbf{ALSO PEPs}.}
	\end{frame}
	
	% ================================================================
	% Q25: DOMESTIC PEPs
	% ================================================================
	\begin{frame}
		\frametitle{\fontsize{12}{14}\selectfont \textbf{Q25: Domestic PEPs — Risk Level}}
		\begin{block}{\fontsize{9}{11}\selectfont \textbf{Question}}
			\scriptsize
			A bank onboards a domestic PEP. The compliance team argues domestic PEPs are ``lower risk'' and standard CDD suffices. Is this correct?
			\vspace{0.1cm}
			\begin{itemize}
				\item[\textbf{A.}] Yes — FATF Rec 12 explicitly states domestic PEPs present lower risk.
				\item[\textbf{B.}] No — domestic PEPs require a risk-based assessment; EDD is required where elevated risk is identified.
				\item[\textbf{C.}] Yes — domestic PEPs are not covered under FATF Rec 12 at all.
				\item[\textbf{D.}] No — domestic PEPs always require the same EDD as foreign PEPs.
			\end{itemize}
		\end{block}
		\vspace{0.05cm}
		\answerbox{\large \textbf{B} is correct.}
		\vspace{0.05cm}
		\stepbox{\scriptsize
			\keypoint{Explanation:} FATF Rec 12 allows a \textbf{risk-based approach} for domestic PEPs. Institutions CANNOT simply say ``domestic = standard CDD'' without conducting the risk assessment. EDD is still required when assessment shows elevated risk.
		}
		\vspace{0.03cm}
		\trapbox{\scriptsize Domestic PEPs need \textbf{RISK ASSESSMENT first}, not automatic standard CDD.}
	\end{frame}
	
	% ================================================================
	% Q26: PEP EDD REQUIREMENTS
	% ================================================================
	\begin{frame}
		\frametitle{\fontsize{12}{14}\selectfont \textbf{Q26: PEP EDD — Required Measures under Rec 12}}
		\begin{block}{\fontsize{9}{11}\selectfont \textbf{Question}}
			\scriptsize
			A bank onboards a foreign PEP with \$20M in assets. Under FATF Recommendation 12, which combination of EDD measures is REQUIRED?
			\vspace{0.1cm}
			\begin{itemize}
				\item[\textbf{A.}] Source of funds verification only — wealth verification is optional.
				\item[\textbf{B.}] Senior management approval, source of WEALTH AND source of FUNDS verification, and enhanced ongoing monitoring.
				\item[\textbf{C.}] Standard CDD plus an annual review.
				\item[\textbf{D.}] Outsource EDD to a third-party screening provider and rely on their certification.
			\end{itemize}
		\end{block}
		\vspace{0.05cm}
		\answerbox{\large \textbf{B} is correct.}
		\vspace{0.05cm}
		\stepbox{\scriptsize
			\keypoint{Explanation:} FATF Recommendation 12 mandates: (1) \textbf{Senior management APPROVAL}; (2) \textbf{SOURCE OF WEALTH} AND \textbf{SOURCE OF FUNDS} — \textit{both required}; (3) \textbf{Enhanced ongoing monitoring}. Third-party screening does NOT discharge the institution's responsibility.
		}
		\vspace{0.03cm}
		\trapbox{\scriptsize \textbf{Source of Wealth} $\neq$ \textbf{Source of Funds}. FATF Rec 12 requires \textbf{BOTH}.}
	\end{frame}
	
	% ================================================================
	% SECTION 14: BENEFICIAL OWNERSHIP
	% ================================================================
	\begin{frame}
		\frametitle{\fontsize{13}{15}\selectfont \textbf{Section 14: Beneficial Ownership}}
		\vspace{0.3cm}
		\begin{center}
			{\fontsize{12}{14}\selectfont \textbf{FATF Recommendation 24 — Key Definition}}
		\end{center}
		\vspace{0.2cm}
		\begin{block}{}
			\centering
			\textbf{Beneficial Owner} = The natural person(s) who ultimately OWN or CONTROL a legal entity.
			\vspace{0.1cm}
			\begin{itemize}
				\item 25\% is a \textbf{starting point}, NOT a safe harbor
				\item Control through other means triggers identification regardless of ownership percentage
				\item Fallback: if no natural person meets the threshold, the senior managing official serves as fallback
			\end{itemize}
		\end{block}
	\end{frame}
	
	% ================================================================
	% Q27: BENEFICIAL OWNERSHIP
	% ================================================================
	\begin{frame}
		\frametitle{\fontsize{12}{14}\selectfont \textbf{Q27: Beneficial Ownership — FATF Recommendation 24}}
		\begin{block}{\fontsize{9}{11}\selectfont \textbf{Question}}
			\scriptsize
			Under FATF Recommendation 24, what defines the ``beneficial owner'' of a company?
			\vspace{0.1cm}
			\begin{itemize}
				\item[\textbf{A.}] The registered director of the company.
				\item[\textbf{B.}] Any person with a legal contract with the company.
				\item[\textbf{C.}] The natural person(s) who ultimately OWN or CONTROL the legal entity (commonly using 25\% as an initial threshold).
				\item[\textbf{D.}] The CEO regardless of equity ownership percentage.
			\end{itemize}
		\end{block}
		\vspace{0.05cm}
		\answerbox{\large \textbf{C} is correct.}
		\vspace{0.05cm}
		\stepbox{\scriptsize
			\keypoint{Explanation:} \textbf{Beneficial owner} = natural person who ultimately owns or controls. Key elements: (1) \textbf{Natural person}; (2) \textbf{Ultimate} — through all layers; (3) \textbf{Control} — through ownership, voting rights, board control; (4) \textbf{25\% threshold} is a starting point only.
		}
		\vspace{0.03cm}
		\trapbox{\scriptsize 25\% is a \textbf{starting point}, not a safe harbor. Control through other means triggers identification.}
	\end{frame}
	
	% ================================================================
	% SECTION 15: CDD
	% ================================================================
	\begin{frame}
		\frametitle{\fontsize{13}{15}\selectfont \textbf{Section 15: Customer Due Diligence (CDD)}}
		\vspace{0.3cm}
		\begin{center}
			{\fontsize{12}{14}\selectfont \textbf{Ongoing Monitoring \& Simplified Due Diligence}}
		\end{center}
		\vspace{0.2cm}
		\begin{block}{\textbf{Ongoing Monitoring}}
			\scriptsize
			Compare ACTUAL transaction activity against the \textbf{EXPECTED PROFILE} established during onboarding. Significant deviations trigger review.
		\end{block}
		\vspace{0.1cm}
		\begin{block}{\textbf{Simplified Due Diligence (SDD)}}
			\scriptsize
			Permitted only when risk is \textbf{LOW} — based on NRA + institutional risk assessment. SDD does NOT mean zero CDD — core identification is still required. \textbf{Suspicion overrides SDD automatically.}
		\end{block}
	\end{frame}
	
	% ================================================================
	% Q28: ONGOING MONITORING
	% ================================================================
	\begin{frame}
		\frametitle{\fontsize{12}{14}\selectfont \textbf{Q28: CDD — Ongoing Monitoring}}
		\begin{block}{\fontsize{9}{11}\selectfont \textbf{Question}}
			\scriptsize
			Which transaction pattern should MOST URGENTLY trigger a CDD review under ongoing monitoring obligations?
			\vspace{0.1cm}
			\begin{itemize}
				\item[\textbf{A.}] A business account receiving monthly revenue consistent with declared activity.
				\item[\textbf{B.}] A personal account receiving a one-time inheritance consistent with prior disclosures.
				\item[\textbf{C.}] A business declared as a small retail shop suddenly receiving \$500K in international wire transfers from high-risk jurisdictions.
				\item[\textbf{D.}] A personal savings account with monthly direct deposit from the same employer for 3 years.
			\end{itemize}
		\end{block}
		\vspace{0.05cm}
		\answerbox{\large \textbf{C} is correct.}
		\vspace{0.05cm}
		\stepbox{\scriptsize
			\keypoint{Explanation:} Ongoing monitoring requires comparing ACTUAL activity against the \textbf{EXPECTED PROFILE}. The highest alert trigger: transactions INCONSISTENT WITH THE CUSTOMER'S DECLARED BUSINESS PURPOSE.
		}
		\vspace{0.03cm}
		\trapbox{\scriptsize Ongoing monitoring = compare actual activity to \textbf{expected profile}.}
	\end{frame}
	
	% ================================================================
	% Q29: SIMPLIFIED DUE DILIGENCE
	% ================================================================
	\begin{frame}
		\frametitle{\fontsize{12}{14}\selectfont \textbf{Q29: Simplified Due Diligence (SDD) — When Permitted}}
		\begin{block}{\fontsize{9}{11}\selectfont \textbf{Question}}
			\scriptsize
			Under FATF Recommendation 10, when is Simplified Due Diligence (SDD) permissible?
			\vspace{0.1cm}
			\begin{itemize}
				\item[\textbf{A.}] Always — SDD is the default for retail customers.
				\item[\textbf{B.}] Never — FATF requires full CDD for all customers at all times.
				\item[\textbf{C.}] Only when the NRA, institutional risk assessment, and customer/product analysis demonstrate the ML/TF risk is LOW.
				\item[\textbf{D.}] Only for government entities and listed public companies.
			\end{itemize}
		\end{block}
		\vspace{0.05cm}
		\answerbox{\large \textbf{C} is correct.}
		\vspace{0.05cm}
		\stepbox{\scriptsize
			\keypoint{Explanation:} SDD is permitted only when risk is \textbf{LOW}. SDD does NOT mean zero CDD — core identification is still required. \textbf{Suspicion overrides SDD automatically.}
		}
		\vspace{0.03cm}
		\trapbox{\scriptsize SDD $\neq$ no CDD. Suspicion \textbf{overrides SDD automatically}.}
	\end{frame}
	
	% ================================================================
	% SECTION 16: MSB
	% ================================================================
	\begin{frame}
		\frametitle{\fontsize{13}{15}\selectfont \textbf{Section 16: Money Services Businesses (MSBs)}}
		\vspace{0.3cm}
		\begin{center}
			{\fontsize{12}{14}\selectfont \textbf{Agent Network Liability — Principal Responsibility}}
		\end{center}
		\vspace{0.2cm}
		\recall{\large \textbf{Principal MSBs are responsible for their agents' AML compliance.}}
		\vspace{0.1cm}
		\begin{block}{}
			\scriptsize
			Under US BSA rules (31 CFR 1022.380):
			\begin{itemize}
				\item Principal registers with FinCEN; agents do not separately register
				\item Principal must maintain a current list of agents
				\item Principal's AML program must cover agent activity
				\item Agent liability = principal liability
			\end{itemize}
			\vspace{0.1cm}
			\textbf{Example:} Western Union enforcement action (2017, \$586M settlement).
		\end{block}
	\end{frame}
	
	% ================================================================
	% Q30: MSB AGENT NETWORK
	% ================================================================
	\begin{frame}
		\frametitle{\fontsize{12}{14}\selectfont \textbf{Q30: MSB Agent Networks — Principal Liability}}
		\begin{block}{\fontsize{9}{11}\selectfont \textbf{Question}}
			\scriptsize
			A large US money transmitter has 500 agents across high-risk jurisdictions with no AML oversight of agent activity. What is the PRINCIPAL regulatory risk?
			\vspace{0.1cm}
			\begin{itemize}
				\item[\textbf{A.}] Currency exchange rate risk from international transactions.
				\item[\textbf{B.}] The principal transmitter bears responsibility for AML compliance across its agent network.
				\item[\textbf{C.}] Agents operate under local rules only; principal has no responsibility.
				\item[\textbf{D.}] Reputation risk only — agent failures create no legal liability.
			\end{itemize}
		\end{block}
		\vspace{0.05cm}
		\answerbox{\large \textbf{B} is correct.}
		\vspace{0.05cm}
		\stepbox{\scriptsize
			\keypoint{Explanation:} Under US BSA rules, a \textbf{principal money transmitter is responsible for AML compliance of its AGENTS}. If an agent facilitates ML, the PRINCIPAL can face regulatory action. The Western Union enforcement action (2017, \$586M) illustrates this.
		}
		\vspace{0.03cm}
		\trapbox{\scriptsize Principal MSBs are \textbf{responsible for their agents'} AML compliance.}
	\end{frame}
	
	% ================================================================
	% SECTION 17: PSP
	% ================================================================
	\begin{frame}
		\frametitle{\fontsize{13}{15}\selectfont \textbf{Section 17: Payment Service Providers (PSPs)}}
		\vspace{0.3cm}
		\begin{center}
			{\fontsize{12}{14}\selectfont \textbf{Transaction Laundering — The Key Risk}}
		\end{center}
		\vspace{0.2cm}
		\begin{block}{}
			\centering
			\textbf{Transaction Laundering:} Using fake/front merchant accounts to process payments for illicit goods/services through a legitimate PSP.
			\vspace{0.1cm}
			\textbf{Red Flags:}
			\begin{itemize}
				\item Minimal real inventory but high transaction volume
				\item Multiple accounts with inconsistent beneficiaries
				\item Refund/chargeback patterns inconsistent with legitimate retail
			\end{itemize}
		\end{block}
	\end{frame}
	
	% ================================================================
	% Q31: PSP RISKS
	% ================================================================
	\begin{frame}
		\frametitle{\fontsize{12}{14}\selectfont \textbf{Q31: PSP / E-Commerce — Transaction Laundering}}
		\begin{block}{\fontsize{9}{11}\selectfont \textbf{Question}}
			\scriptsize
			An online marketplace detects sellers with low inventory but processing thousands of small transactions, using multiple account names with different payment destinations. What ML technique does this represent?
			\vspace{0.1cm}
			\begin{itemize}
				\item[\textbf{A.}] Legitimate high-volume retail — no concern.
				\item[\textbf{B.}] Transaction laundering — using fake/front merchant accounts to process payments for illicit goods/services.
				\item[\textbf{C.}] Arbitrage trading — legal profit from price differentials.
				\item[\textbf{D.}] Chargeback fraud — disputing legitimate transactions.
			\end{itemize}
		\end{block}
		\vspace{0.05cm}
		\answerbox{\large \textbf{B} is correct.}
		\vspace{0.05cm}
		\stepbox{\scriptsize
			\keypoint{Explanation:} \textbf{Transaction laundering}: illicit actors use legitimate merchant accounts to process payments for illegal goods/services. Pattern: minimal real inventory but high transaction volume; multiple accounts with inconsistent beneficiaries.
		}
		\vspace{0.03cm}
		\trapbox{\scriptsize Transaction laundering = illicit proceeds hidden as e-commerce revenue.}
	\end{frame}
	
	% ================================================================
	% Q32: E-COMMERCE MITIGATION
	% ================================================================
	\begin{frame}
		\frametitle{\fontsize{12}{14}\selectfont \textbf{Q32: E-Commerce ML — Mitigation Controls}}
		\begin{block}{\fontsize{9}{11}\selectfont \textbf{Question}}
			\scriptsize
			An e-commerce platform wants to reduce ML risk while maintaining business growth. Which combination of controls is MOST EFFECTIVE?
			\vspace{0.1cm}
			\begin{itemize}
				\item[\textbf{A.}] Require all sellers to use cryptocurrency to enable blockchain tracing.
				\item[\textbf{B.}] Tiered seller onboarding, transaction monitoring, sanctions screening, and STR/SAR filing.
				\item[\textbf{C.}] Limit the platform to domestic sellers only.
				\item[\textbf{D.}] Only apply controls to sellers transacting over \$10,000 per month.
			\end{itemize}
		\end{block}
		\vspace{0.05cm}
		\answerbox{\large \textbf{B} is correct.}
		\vspace{0.05cm}
		\stepbox{\scriptsize
			\keypoint{Explanation:} Effective e-commerce AML controls: (1) \textbf{Tiered onboarding} — scale CDD to risk; (2) \textbf{Transaction monitoring}; (3) \textbf{Sanctions screening}; (4) \textbf{STR/SAR filing}. Cryptocurrency does not resolve ML risk.
		}
		\vspace{0.03cm}
		\trapbox{\scriptsize Cryptocurrency does not resolve ML risk — it creates additional VASP obligations.}
	\end{frame}
	
	% ================================================================
	% SECTION 18: INSURANCE ML
	% ================================================================
	\begin{frame}
		\frametitle{\fontsize{13}{15}\selectfont \textbf{Section 18: Life Insurance ML Risks}}
		\vspace{0.3cm}
		\begin{center}
			{\fontsize{12}{14}\selectfont \textbf{Placement and Integration Through Insurance}}
		\end{center}
		\vspace{0.2cm}
		\begin{block}{}
			\centering
			\textbf{Highest-Risk Combination:}
			\vspace{0.1cm}
			\begin{enumerate}
				\item \textbf{Single large premium} (lump sum = placement)
				\item \textbf{Cash payment} (avoids bank KYC)
				\item \textbf{Early surrender} immediately after issue (integration)
			\end{enumerate}
			\vspace{0.1cm}
			The penalty for early surrender is the ``laundering fee.''
		\end{block}
	\end{frame}
	
	% ================================================================
	% Q33: LIFE INSURANCE ML
	% ================================================================
	\begin{frame}
		\frametitle{\fontsize{12}{14}\selectfont \textbf{Q33: Life Insurance — ML Red Flags}}
		\begin{block}{\fontsize{9}{11}\selectfont \textbf{Question}}
			\scriptsize
			Which combination makes a life insurance policy HIGHEST risk for money laundering?
			\vspace{0.1cm}
			\begin{itemize}
				\item[\textbf{A.}] A 30-year term life policy with monthly premiums from a salary account.
				\item[\textbf{B.}] A single-premium life policy (\$500K+) purchased with cash, where the policyholder requests early surrender immediately after issue.
				\item[\textbf{C.}] A group life policy purchased by an employer for 50 employees.
				\item[\textbf{D.}] A universal life policy with automatic premium waiver.
			\end{itemize}
		\end{block}
		\vspace{0.05cm}
		\answerbox{\large \textbf{B} is correct.}
		\vspace{0.05cm}
		\stepbox{\scriptsize
			\keypoint{Explanation:} Highest-risk combination: (1) \textbf{Single large premium}; (2) \textbf{Cash payment}; (3) \textbf{Early surrender} — converts illicit cash $\rightarrow$ insurance premium $\rightarrow$ insurance check (appears legitimate).
		}
		\vspace{0.03cm}
		\trapbox{\scriptsize Single large premium + early surrender = classic \textbf{placement/integration} cycle.}
	\end{frame}
	
	% ================================================================
	% SECTION 19: CASINO ML
	% ================================================================
	\begin{frame}
		\frametitle{\fontsize{13}{15}\selectfont \textbf{Section 19: Casino ML Risks}}
		\vspace{0.3cm}
		\begin{center}
			{\fontsize{12}{14}\selectfont \textbf{The Classic Casino ML Pattern}}
		\end{center}
		\vspace{0.2cm}
		\begin{block}{}
			\centering
			\textbf{Placement:} Cash $\rightarrow$ casino chips \\
			\textbf{Minimal gambling:} The ``loss'' is the laundering service fee \\
			\textbf{Integration:} Casino check with ``gambling winnings'' paper trail
			\vspace{0.1cm}
			\recall{Casinos ARE covered by BSA SAR and CTR obligations (31 CFR 1021).}
		\end{block}
	\end{frame}
	
	% ================================================================
	% Q34: CASINO ML
	% ================================================================
	\begin{frame}
		\frametitle{\fontsize{12}{14}\selectfont \textbf{Q34: Casino ML — Chip Placement and Integration}}
		\begin{block}{\fontsize{9}{11}\selectfont \textbf{Question}}
			\scriptsize
			A casino high-roller buys \$200,000 in chips with cash, gambles minimally, then cashes out \$190,000 as a casino check. What ML stage does the \$10,000 ``loss'' represent, and why is the casino check significant?
			\vspace{0.1cm}
			\begin{itemize}
				\item[\textbf{A.}] Layering only — the casino check obscures the cash origin.
				\item[\textbf{B.}] Placement (chip purchase with cash) followed by integration (casino check creates a ``gambling winnings'' paper trail).
				\item[\textbf{C.}] Structuring — \$10,000 below CTR threshold.
				\item[\textbf{D.}] Tax evasion — gambling losses reduce taxable income.
			\end{itemize}
		\end{block}
		\vspace{0.05cm}
		\answerbox{\large \textbf{B} is correct.}
		\vspace{0.05cm}
		\stepbox{\scriptsize
			\keypoint{Explanation:} Classic casino ML: (1) \textbf{Placement} — cash to chips; (2) Minimal gambling (\$10,000 loss = laundering fee); (3) \textbf{Integration} — casino check with ``gambling winnings'' paper trail. The \$200K chip purchase triggers a \textbf{CTR}.
		}
		\vspace{0.03cm}
		\trapbox{\scriptsize Casino check = \textbf{integration tool}. The \$200K chip purchase requires a \textbf{CTR}.}
	\end{frame}
	
	% ================================================================
	% SECTION 20: REAL ESTATE ML
	% ================================================================
	\begin{frame}
		\frametitle{\fontsize{13}{15}\selectfont \textbf{Section 20: Real Estate ML Risks}}
		\vspace{0.3cm}
		\begin{center}
			{\fontsize{12}{14}\selectfont \textbf{Real Estate as an Integration Vehicle}}
		\end{center}
		\vspace{0.2cm}
		\begin{block}{}
			\centering
			\textbf{Highest-Risk Combination:}
			\vspace{0.1cm}
			\begin{enumerate}
				\item \textbf{Shell company} — conceals beneficial owner
				\item \textbf{All-cash purchase} — bypasses bank mortgage KYC
				\item \textbf{Unknown beneficial owner} — critical AML gap
				\item \textbf{Foreign buyer} — potential cross-border ML
			\end{enumerate}
		\end{block}
	\end{frame}
	
	% ================================================================
	% Q35: REAL ESTATE ML
	% ================================================================
	\begin{frame}
		\frametitle{\fontsize{12}{14}\selectfont \textbf{Q35: Real Estate — ML Integration}}
		\begin{block}{\fontsize{9}{11}\selectfont \textbf{Question}}
			\scriptsize
			A foreign shell company purchases a \$3M residential property with all-cash (no mortgage). The beneficial owner is unknown. Which red flag combination is MOST significant?
			\vspace{0.1cm}
			\begin{itemize}
				\item[\textbf{A.}] Residential property — commercial real estate is higher risk.
				\item[\textbf{B.}] Shell company + all-cash purchase + unknown beneficial owner + no mortgage = classic real estate ML integration pattern.
				\item[\textbf{C.}] Foreign buyer only — domestic purchases present no ML risk.
				\item[\textbf{D.}] Amount is below significant thresholds; no reporting required.
			\end{itemize}
		\end{block}
		\vspace{0.05cm}
		\answerbox{\large \textbf{B} is correct.}
		\vspace{0.05cm}
		\stepbox{\scriptsize
			\keypoint{Explanation:} Highest-risk combination: (1) \textbf{Shell company}; (2) \textbf{All-cash purchase} (bypasses bank KYC); (3) \textbf{Unknown beneficial owner}. FinCEN GTOs require beneficial owner identification for all-cash shell company purchases.
		}
		\vspace{0.03cm}
		\trapbox{\scriptsize Cash purchase = \textbf{no bank KYC layer}. Shell company + cash = highest risk.}
	\end{frame}
	
	% ================================================================
	% SECTION 21: ART ML
	% ================================================================
	\begin{frame}
		\frametitle{\fontsize{13}{15}\selectfont \textbf{Section 21: Art and Antiquities ML Risks}}
		\vspace{0.3cm}
		\begin{center}
			{\fontsize{12}{14}\selectfont \textbf{Why Art Is Attractive for ML}}
		\end{center}
		\vspace{0.2cm}
		\begin{block}{}
			\scriptsize
			\begin{itemize}
				\item \textbf{Subjective valuation} — no objective market price, enabling manipulation
				\item \textbf{Private sales} — no exchange, no public price discovery, anonymous parties
				\item \textbf{Portability} — crosses borders with limited inspection
				\item \textbf{Over/under-invoicing} — transfer value between parties
			\end{itemize}
			\vspace{0.1cm}
			\textbf{US AMLA 2020:} Expanded BSA to cover art dealers for transactions over \$10,000.
		\end{block}
	\end{frame}
	
	% ================================================================
	% Q36: ART ML
	% ================================================================
	\begin{frame}
		\frametitle{\fontsize{12}{14}\selectfont \textbf{Q36: Art and Antiquities — ML Risks}}
		\begin{block}{\fontsize{9}{11}\selectfont \textbf{Question}}
			\scriptsize
			A high-value artwork is sold through a private sale with no public record, at a price far above any prior valuation. What ML risk does this represent?
			\vspace{0.1cm}
			\begin{itemize}
				\item[\textbf{A.}] No ML risk — art sales are private transactions not subject to AML.
				\item[\textbf{B.}] Integration via subjective valuation: art prices are subjective, enabling over-valuation to justify moving large sums.
				\item[\textbf{C.}] Tax avoidance only — art sales are not covered by FATF.
				\item[\textbf{D.}] Placement only — art purchases introduce new illicit funds.
			\end{itemize}
		\end{block}
		\vspace{0.05cm}
		\answerbox{\large \textbf{B} is correct.}
		\vspace{0.05cm}
		\stepbox{\scriptsize
			\keypoint{Explanation:} Art is attractive for ML integration due to: (1) \textbf{Subjective valuation}; (2) \textbf{Private sales}; (3) \textbf{Portability}; (4) \textbf{Over/under-invoicing}. The US AMLA 2020 expanded BSA to cover art dealers.
		}
		\vspace{0.03cm}
		\trapbox{\scriptsize Subjective valuation = core art ML risk. Art dealers are now \textbf{covered entities}.}
	\end{frame}
	
	% ================================================================
	% SECTION 22: VASPs
	% ================================================================
	\begin{frame}
		\frametitle{\fontsize{13}{15}\selectfont \textbf{Section 22: Virtual Asset Service Providers (VASPs)}}
		\vspace{0.3cm}
		\begin{center}
			{\fontsize{12}{14}\selectfont \textbf{FATF Recommendations 15 \& 16}}
		\end{center}
		\vspace{0.2cm}
		\begin{block}{}
			\centering
			\textbf{Rec 15:} Foundational framework — VASPs must be subject to AML/CFT regulation. \\
			\textbf{Rec 16 (Travel Rule):} Requires originator and beneficiary information for VA transfers $\geq$\$1,000.
			\vspace{0.1cm}
			\textbf{Privacy Coins (Monero, Zcash):} Use cryptographic techniques making blockchain analysis ineffective.
		\end{block}
	\end{frame}
	
	% ================================================================
	% Q37: VASP — REC 15
	% ================================================================
	\begin{frame}
		\frametitle{\fontsize{12}{14}\selectfont \textbf{Q37: VASPs — FATF Recommendation 15}}
		\begin{block}{\fontsize{9}{11}\selectfont \textbf{Question}}
			\scriptsize
			Under FATF standards, which Recommendation is the PRIMARY framework governing Virtual Asset Service Providers (VASPs)?
			\vspace{0.1cm}
			\begin{itemize}
				\item[\textbf{A.}] Recommendation 13 — Correspondent Banking.
				\item[\textbf{B.}] Recommendation 16 — Wire Transfers (Travel Rule).
				\item[\textbf{C.}] Recommendation 15 — New Technologies, updated in 2019 to include virtual assets.
				\item[\textbf{D.}] Recommendation 22 — DNFBPs.
			\end{itemize}
		\end{block}
		\vspace{0.05cm}
		\answerbox{\large \textbf{C} is correct.}
		\vspace{0.05cm}
		\stepbox{\scriptsize
			\keypoint{Explanation:} FATF \textbf{Recommendation 15} covers New Technologies and requires VASPs to be subject to AML/CFT regulation. \textbf{Recommendation 16} (Travel Rule) ALSO applies to VASPs for virtual asset transfers — but Rec 15 is the FOUNDATIONAL framework.
		}
		\vspace{0.03cm}
		\trapbox{\scriptsize Primary VASP framework = \textbf{Rec 15}. Travel Rule = \textbf{Rec 16}.}
	\end{frame}
	
	% ================================================================
	% Q38: TRAVEL RULE
	% ================================================================
	\begin{frame}
		\frametitle{\fontsize{12}{14}\selectfont \textbf{Q38: VASP Travel Rule — Required Information}}
		\begin{block}{\fontsize{9}{11}\selectfont \textbf{Question}}
			\scriptsize
			Under the FATF Travel Rule (Recommendation 16) extended to VASPs, which information must accompany a virtual asset transfer of \$1,000+?
			\vspace{0.1cm}
			\begin{itemize}
				\item[\textbf{A.}] Only the originator's wallet address — beneficiary info is optional.
				\item[\textbf{B.}] Originator: full name, wallet address, AND one of (physical address/national ID/customer number/DOB); Beneficiary: name and wallet address.
				\item[\textbf{C.}] The transaction hash and block confirmation number only.
				\item[\textbf{D.}] AML risk score for the wallet, but no personal identifying information.
			\end{itemize}
		\end{block}
		\vspace{0.05cm}
		\answerbox{\large \textbf{B} is correct.}
		\vspace{0.05cm}
		\stepbox{\scriptsize
			\keypoint{Explanation:} FATF Rec 16 extended to VASPs requires: \textbf{ORIGINATOR}: full name, wallet address, AND at least one identifier; \textbf{BENEFICIARY}: name and wallet address. \textbf{Privacy coins and unhosted wallets} are major compliance gaps.
		}
		\vspace{0.03cm}
		\trapbox{\scriptsize Both \textbf{originator AND beneficiary} information required. Wallet address alone is insufficient.}
	\end{frame}
	
	% ================================================================
	% Q39: PRIVACY COINS
	% ================================================================
	\begin{frame}
		\frametitle{\fontsize{12}{14}\selectfont \textbf{Q39: Privacy Coins vs Standard Cryptocurrency}}
		\begin{block}{\fontsize{9}{11}\selectfont \textbf{Question}}
			\scriptsize
			Why do privacy coins (e.g., Monero, Zcash) present GREATER AML risk than standard cryptocurrencies like Bitcoin?
			\vspace{0.1cm}
			\begin{itemize}
				\item[\textbf{A.}] Privacy coins are faster, allowing criminals to move more funds.
				\item[\textbf{B.}] Privacy coins obfuscate sender, receiver, and transaction amounts using cryptographic techniques, making blockchain analysis largely ineffective.
				\item[\textbf{C.}] Privacy coins are not covered under FATF Recommendation 15.
				\item[\textbf{D.}] Privacy coins require more computing power, making them harder to confiscate.
			\end{itemize}
		\end{block}
		\vspace{0.05cm}
		\answerbox{\large \textbf{B} is correct.}
		\vspace{0.05cm}
		\stepbox{\scriptsize
			\keypoint{Explanation:} \textbf{Bitcoin} = pseudonymous, NOT anonymous — blockchain analytics can trace transactions. \textbf{Privacy coins}: Monero uses ring signatures, stealth addresses, RingCT — making tracing near-impossible. Zcash uses zero-knowledge proofs (zk-SNARKs).
		}
		\vspace{0.03cm}
		\trapbox{\scriptsize Bitcoin = \textbf{pseudonymous} (traceable). Privacy coins = \textbf{cryptographically anonymous}.}
	\end{frame}
	
	% ================================================================
	% SECTION 23: SECURITIES SECTOR
	% ================================================================
	\begin{frame}
		\frametitle{\fontsize{13}{15}\selectfont \textbf{Section 23: Securities Sector — Wash Trading}}
		\vspace{0.3cm}
		\begin{center}
			{\fontsize{12}{14}\selectfont \textbf{Market Manipulation and ML Integration}}
		\end{center}
		\vspace{0.2cm}
		\begin{block}{}
			\centering
			\textbf{Wash Trading:} Executing buy and sell orders for the same security between related accounts to create the APPEARANCE of trading activity.
			\vspace{0.1cm}
			\textbf{In ML context:}
			\begin{itemize}
				\item Illicit funds enter as ``investment capital''
				\item Round-trip trades generate ``trading profit'' documentation
				\item Proceeds withdrawn as ``investment returns'' — legitimized (integration)
			\end{itemize}
		\end{block}
	\end{frame}
	
	% ================================================================
	% Q40: SECURITIES WASH TRADING
	% ================================================================
	\begin{frame}
		\frametitle{\fontsize{12}{14}\selectfont \textbf{Q40: Securities Sector — Wash Trading}}
		\begin{block}{\fontsize{9}{11}\selectfont \textbf{Question}}
			\scriptsize
			A securities firm identifies a customer conducting rapid round-trip trades — buying and selling the same securities between related accounts within hours, with no net economic gain. What does this indicate?
			\vspace{0.1cm}
			\begin{itemize}
				\item[\textbf{A.}] High-frequency trading — a legal and common investment strategy.
				\item[\textbf{B.}] Wash trading — creating artificial trading activity to generate apparent profits for ML integration or market manipulation.
				\item[\textbf{C.}] Dollar-cost averaging — a conservative investment technique.
				\item[\textbf{D.}] Short selling — profiting from declining share prices.
			\end{itemize}
		\end{block}
		\vspace{0.05cm}
		\answerbox{\large \textbf{B} is correct.}
		\vspace{0.05cm}
		\stepbox{\scriptsize
			\keypoint{Explanation:} \textbf{Wash trading} is ML integration AND market manipulation. Securities firms must file both \textbf{SARs} (ML component) AND \textbf{market abuse reports} (manipulation component).
		}
		\vspace{0.03cm}
		\trapbox{\scriptsize Round-trips with no economic gain = \textbf{wash trading} = ML integration AND market manipulation.}
	\end{frame}
	
	% ================================================================
	% SECTION 24: GATEKEEPERS
	% ================================================================
	\begin{frame}
		\frametitle{\fontsize{13}{15}\selectfont \textbf{Section 24: Gatekeepers — Lawyers \& TCSPs}}
		\vspace{0.3cm}
		\begin{center}
			{\fontsize{12}{14}\selectfont \textbf{FATF Recommendation 23 — DNFBP Obligations}}
		\end{center}
		\vspace{0.2cm}
		\begin{block}{}
			\scriptsize
			\textbf{Lawyers:}
			\begin{itemize}
				\item Can be criminally prosecuted for facilitating ML with knowledge of the predicate offense
				\item Attorney-client privilege protects \textbf{legal advice}, NOT \textbf{transactional facilitation}
				\item Must file STRs for designated activities (real estate, company formation, fund management)
			\end{itemize}
			\vspace{0.1cm}
			\textbf{TCSPs:}
			\begin{itemize}
				\item Nominee directors/shareholders obscure beneficial ownership
				\item Must identify the TRUE beneficial owner — not just the nominee client
				\item FATF Recs 24/25 require accurate beneficial ownership information
			\end{itemize}
		\end{block}
	\end{frame}
	
	% ================================================================
	% Q41: LAWYER LIABILITY
	% ================================================================
	\begin{frame}
		\frametitle{\fontsize{12}{14}\selectfont \textbf{Q41: Gatekeepers — Lawyer Liability for ML}}
		\begin{block}{\fontsize{9}{11}\selectfont \textbf{Question}}
			\scriptsize
			A lawyer helps a client incorporate a company, open bank accounts, and purchase real estate, knowing the client's funds derive from drug trafficking. Under what principle can the lawyer face criminal liability?
			\vspace{0.1cm}
			\begin{itemize}
				\item[\textbf{A.}] Lawyers are fully exempt from ML liability due to attorney-client privilege.
				\item[\textbf{B.}] Only civil sanctions apply to lawyers — no criminal prosecution possible.
				\item[\textbf{C.}] The lawyer can be criminally prosecuted for money laundering by facilitating ML transactions with knowledge of the predicate offense.
				\item[\textbf{D.}] Liability only attaches if the lawyer personally transferred the funds.
			\end{itemize}
		\end{block}
		\vspace{0.05cm}
		\answerbox{\large \textbf{C} is correct.}
		\vspace{0.05cm}
		\stepbox{\scriptsize
			\keypoint{Explanation:} \textbf{Attorney-client privilege} protects CONFIDENTIAL COMMUNICATIONS for legal advice — it does NOT protect a lawyer from criminal prosecution for PARTICIPATING IN or FACILITATING financial crime with knowledge.
		}
		\vspace{0.03cm}
		\trapbox{\scriptsize Privilege protects \textbf{legal advice}. It does NOT immunize \textbf{transactional participation} in ML.}
	\end{frame}
	
	% ================================================================
	% Q42: TCSP RISKS
	% ================================================================
	\begin{frame}
		\frametitle{\fontsize{12}{14}\selectfont \textbf{Q42: TCSPs — Nominee Director Risk}}
		\begin{block}{\fontsize{9}{11}\selectfont \textbf{Question}}
			\scriptsize
			A TCSP provides nominee director and shareholder services without verifying the ultimate beneficial owner. What is the PRIMARY ML risk?
			\vspace{0.1cm}
			\begin{itemize}
				\item[\textbf{A.}] The TCSP creates currency exchange risk.
				\item[\textbf{B.}] Nominee arrangements obscure true beneficial ownership, enabling criminals to control companies while remaining invisible.
				\item[\textbf{C.}] TCSPs are exempt from beneficial ownership requirements.
				\item[\textbf{D.}] The risk is limited to the jurisdiction where the trust is formed.
			\end{itemize}
		\end{block}
		\vspace{0.05cm}
		\answerbox{\large \textbf{B} is correct.}
		\vspace{0.05cm}
		\stepbox{\scriptsize
			\keypoint{Explanation:} TCSPs are high-risk because nominee directors/shareholders appear in public records \textbf{in place of the true beneficial owner}. TCSPs must apply CDD to identify the TRUE beneficial owner.
		}
		\vspace{0.03cm}
		\trapbox{\scriptsize Nominees obscure beneficial ownership. TCSPs must look \textbf{THROUGH} nominees.}
	\end{frame}
	
	% ================================================================
	% SECTION 25: HUMAN TRAFFICKING
	% ================================================================
	\begin{frame}
		\frametitle{\fontsize{13}{15}\selectfont \textbf{Section 25: Human Trafficking — Financial Indicators}}
		\vspace{0.3cm}
		\begin{center}
			{\fontsize{12}{14}\selectfont \textbf{Detecting Trafficking Through Financial Patterns}}
		\end{center}
		\vspace{0.2cm}
		\begin{block}{}
			\scriptsize
			\textbf{Red Flags:}
			\begin{itemize}
				\item Multiple small \textbf{cash deposits from various locations} (victim payments collected by traffickers)
				\item \textbf{Rapid wire transfers abroad} (sending proceeds to traffickers in home country)
				\item \textbf{Multiple rental payments} (controlling multiple locations where victims are housed/controlled)
				\item Third-party payments for accommodation, transportation
				\item Bulk purchases of food, phones, transportation (logistics for victims)
			\end{itemize}
			\vspace{0.1cm}
			\keypoint{Financial institutions are often the ONLY institutional contact point able to detect trafficking through financial patterns.}
		\end{block}
	\end{frame}
	
	% ================================================================
	% Q43: HUMAN TRAFFICKING
	% ================================================================
	\begin{frame}
		\frametitle{\fontsize{12}{14}\selectfont \textbf{Q43: Human Trafficking — Financial Indicators}}
		\begin{block}{\fontsize{9}{11}\selectfont \textbf{Question}}
			\scriptsize
			A bank account receives multiple small cash deposits from different locations, followed by rapid wire transfers abroad, with the account holder paying for multiple apartment rentals. What financial crime might this indicate?
			\vspace{0.1cm}
			\begin{itemize}
				\item[\textbf{A.}] Import/export business activity.
				\item[\textbf{B.}] Human trafficking proceeds — cash from victims, rapid repatriation, multiple rental locations.
				\item[\textbf{C.}] Standard remittance activity by a migrant worker.
				\item[\textbf{D.}] A legitimate real estate portfolio.
			\end{itemize}
		\end{block}
		\vspace{0.05cm}
		\answerbox{\large \textbf{B} is correct.}
		\vspace{0.05cm}
		\stepbox{\scriptsize
			\keypoint{Explanation:} Human trafficking financial signature: \textbf{multiple small cash deposits} (victim payments) + \textbf{rapid overseas wires} (repatriation) + \textbf{multiple rental payments} (controlling locations). FATF issued guidance on this in 2018.
		}
		\vspace{0.03cm}
		\trapbox{\scriptsize Human trafficking signature: small cash deposits + rapid overseas wires + multiple rentals.}
	\end{frame}
	
	% ================================================================
	% SECTION 26: PROLIFERATION FINANCING
	% ================================================================
	\begin{frame}
		\frametitle{\fontsize{13}{15}\selectfont \textbf{Section 26: Proliferation Financing}}
		\vspace{0.3cm}
		\begin{center}
			{\fontsize{12}{14}\selectfont \textbf{FATF Recommendation 7 — PF vs TF}}
		\end{center}
		\vspace{0.2cm}
		\begin{columns}[T]
			\begin{column}{0.48\textwidth}
				\begin{block}{\textbf{Proliferation Financing}}
					\scriptsize
					\begin{itemize}
						\item Financing WMD programs
						\item FATF Rec 7 — targeted financial sanctions
						\item UN Security Council Resolutions
						\item Often involves larger transactions for dual-use goods
					\end{itemize}
				\end{block}
			\end{column}
			\begin{column}{0.48\textwidth}
				\begin{block}{\textbf{Terrorist Financing}}
					\scriptsize
					\begin{itemize}
						\item Financing terrorist acts
						\item FATF Rec 5 — criminalization
						\item Can involve small amounts
						\item Funds can be legitimate or illicit
					\end{itemize}
				\end{block}
			\end{column}
		\end{columns}
	\end{frame}
	
	% ================================================================
	% Q44: PROLIFERATION FINANCING
	% ================================================================
	\begin{frame}
		\frametitle{\fontsize{12}{14}\selectfont \textbf{Q44: Proliferation Financing — FATF Rec 7}}
		\begin{block}{\fontsize{9}{11}\selectfont \textbf{Question}}
			\scriptsize
			How does proliferation financing (PF) differ from terrorist financing (TF), and which FATF Recommendation specifically addresses PF?
			\vspace{0.1cm}
			\begin{itemize}
				\item[\textbf{A.}] PF and TF are identical — both covered under FATF Recommendation 5.
				\item[\textbf{B.}] PF funds weapons of mass destruction programs; TF funds terrorist acts. PF is addressed through targeted financial sanctions under FATF Recommendation 7.
				\item[\textbf{C.}] PF applies only to nuclear programs; TF covers all other weapons.
				\item[\textbf{D.}] PF is addressed under Recommendation 15 as a new technology risk.
			\end{itemize}
		\end{block}
		\vspace{0.05cm}
		\answerbox{\large \textbf{B} is correct.}
		\vspace{0.05cm}
		\stepbox{\scriptsize
			\keypoint{Explanation:} \textbf{Proliferation Financing (PF)} = financing WMD programs. \textbf{Terrorist Financing (TF)} = financing terrorist acts. FATF framework: \textbf{Rec 7} requires targeted financial sanctions related to proliferation per UN Security Council Resolutions.
		}
		\vspace{0.03cm}
		\trapbox{\scriptsize PF $\neq$ TF. \textbf{PF = WMD financing $\rightarrow$ FATF Rec 7}. \textbf{TF = terrorist acts $\rightarrow$ FATF Rec 5}.}
	\end{frame}
	
	% ================================================================
	% SECTION 27: FATF MUTUAL EVALUATIONS
	% ================================================================
	\begin{frame}
		\frametitle{\fontsize{13}{15}\selectfont \textbf{Section 27: FATF Mutual Evaluations}}
		\vspace{0.3cm}
		\begin{center}
			{\fontsize{12}{14}\selectfont \textbf{Two Dimensions of Assessment}}
		\end{center}
		\vspace{0.2cm}
		\begin{block}{}
			\scriptsize
			\textbf{1. Technical Compliance:} Do laws/regulations technically meet FATF's 40 Recommendations? (Rated C/LC/PC/NC)
			\vspace{0.1cm}
			\textbf{2. Effectiveness:} Is the system working in practice with expected outcomes? (11 Immediate Outcomes; rated HE/SE/ME/LE)
			\vspace{0.1cm}
			\keypoint{Perfect laws $\neq$ effectiveness.}
		\end{block}
		\vspace{0.1cm}
		\begin{block}{\textbf{Gray List vs Black List}}
			\scriptsize
			\textbf{Gray List (Increased Monitoring):} EDD required per Rec 19 — not blanket termination. \\
			\textbf{Black List (Call for Action):} Countermeasures up to refusing transactions. \\
			\recall{FATF opposes de-risking as inconsistent with the risk-based approach.}
		\end{block}
	\end{frame}
	
	% ================================================================
	% Q45: FATF MUTUAL EVALUATIONS
	% ================================================================
	\begin{frame}
		\frametitle{\fontsize{12}{14}\selectfont \textbf{Q45: FATF Mutual Evaluations — Two Dimensions}}
		\begin{block}{\fontsize{9}{11}\selectfont \textbf{Question}}
			\scriptsize
			In a FATF Mutual Evaluation (4th Round), countries are assessed on two dimensions. What are they, and what does a Gray List placement require financial institutions to do?
			\vspace{0.1cm}
			\begin{itemize}
				\item[\textbf{A.}] GDP size and number of SARs filed; Gray List requires terminating all relationships.
				\item[\textbf{B.}] Technical compliance AND Effectiveness; Gray List requires EDD per Rec 19, not blanket termination.
				\item[\textbf{C.}] Number of ML convictions and FIU budget; Gray List is informational only.
				\item[\textbf{D.}] Treaty ratification and court convictions; Gray List prohibits all transactions.
			\end{itemize}
		\end{block}
		\vspace{0.05cm}
		\answerbox{\large \textbf{B} is correct.}
		\vspace{0.05cm}
		\stepbox{\scriptsize
			\keypoint{Explanation:} FATF 4th Round assesses: (1) \textbf{Technical Compliance}; (2) \textbf{Effectiveness}. \textbf{Gray List}: FATF Rec 19 requires EDD — not blanket termination. \textbf{De-risking} is NOT compliant with FATF's risk-based approach.
		}
		\vspace{0.03cm}
		\trapbox{\scriptsize Two dimensions: \textbf{technical compliance} + \textbf{effectiveness}. Gray List = \textbf{EDD required}.}
	\end{frame}
	
	% ================================================================
	% SECTION 28: GLOBAL AFC STANDARDS & GUIDANCE
	% ================================================================
	\begin{frame}
		\frametitle{\fontsize{13}{15}\selectfont \textbf{Section 28: Global AFC Standards \& Guidance}}
		\vspace{0.3cm}
		\begin{center}
			{\fontsize{12}{14}\selectfont \textbf{Key Organizations \& Their Roles}}
		\end{center}
		\vspace{0.1cm}
		\begin{block}{\textbf{G-20 Anti-Corruption Working Group}}
			\scriptsize
			Coordinates with member states to establish and implement multi-year anti-corruption action plans.
		\end{block}
		\begin{block}{\textbf{FATF Recommendation 18}}
			\scriptsize
			Requires group-wide AML programs to apply to all branches and subsidiaries.
		\end{block}
		\begin{block}{\textbf{Basel AML Index}}
			\scriptsize
			Compiles vulnerability indicators across legal, political, and financial dimensions for country risk assessment.
		\end{block}
		\begin{block}{\textbf{FATF Immediate Outcome 5}}
			\scriptsize
			Assesses whether legal persons and arrangements are not misused for ML/TF — weak when anonymous shell entities make tracing beneficial ownership difficult.
		\end{block}
	\end{frame}
	
	% ================================================================
	% SECTION 29: TECHNOLOGY FOR AFC COMPLIANCE
	% ================================================================
	\begin{frame}
		\frametitle{\fontsize{13}{15}\selectfont \textbf{Section 29: Technology for AFC Compliance}}
		\vspace{0.3cm}
		\begin{center}
			{\fontsize{12}{14}\selectfont \textbf{AI, RPA, and Device Intelligence}}
		\end{center}
		\vspace{0.1cm}
		\begin{block}{\textbf{AI-Based Screening Systems}}
			\scriptsize
			When inconsistent transaction rejections occur across business units, best first step: \textbf{Engage the governance committee to assess model performance}.
		\end{block}
		\begin{block}{\textbf{Device Intelligence Tools}}
			\scriptsize
			Capture: \textbf{IP address and geolocation}, \textbf{Operating system and installed plugins}, \textbf{Browser type and screen resolution}.
		\end{block}
		\begin{block}{\textbf{Robotic Process Automation (RPA)}}
			\scriptsize
			Reduces compliance costs by handling large-volume, repeatable tasks.
		\end{block}
	\end{frame}
	
	% ================================================================
	% SECTION 30: DATA COLLECTION & PREPARATION
	% ================================================================
	\begin{frame}
		\frametitle{\fontsize{13}{15}\selectfont \textbf{Section 30: Data Collection \& Preparation}}
		\vspace{0.3cm}
		\begin{center}
			{\fontsize{12}{14}\selectfont \textbf{Internal vs External Sources, Data Lineage}}
		\end{center}
		\vspace{0.1cm}
		\begin{block}{\textbf{External Sources for Sanctions Screening}}
			\scriptsize
			FATF blacklists, UN Security Council lists, State Department embargo registries.
		\end{block}
		\begin{block}{\textbf{Internal vs External Sources}}
			\scriptsize
			\textbf{Internal:} Observed customer behavior or system-held values. \\
			\textbf{External:} Sanction or adverse media lists from agencies.
		\end{block}
		\begin{block}{\textbf{Data Lineage for Transparency}}
			\scriptsize
			Maintaining history of data transformations AND capturing source of each data field.
		\end{block}
		\begin{block}{\textbf{Clustering in Transaction Monitoring}}
			\scriptsize
			Reveals shared patterns among entities to detect potential financial crime.
		\end{block}
	\end{frame}
	
	% ================================================================
	% SECTION 31: DNFBPs & FREE TRADE ZONES
	% ================================================================
	\begin{frame}
		\frametitle{\fontsize{13}{15}\selectfont \textbf{Section 31: DNFBPs \& Free Trade Zones}}
		\vspace{0.3cm}
		\begin{center}
			{\fontsize{12}{14}\selectfont \textbf{Designated Non-Financial Businesses \& Professions}}
		\end{center}
		\vspace{0.1cm}
		\begin{block}{\textbf{Managing DNFBP Risks}}
			\scriptsize
			Creating a reputational risk or client selection committee, and using standardized sector-specific questionnaires.
		\end{block}
		\begin{block}{\textbf{Free Trade Zone Risk}}
			\scriptsize
			Absence of standardized inspection protocols for transshipped goods significantly increases ML risk.
		\end{block}
		\begin{block}{\textbf{TCSP Client Acceptance}}
			\scriptsize
			Should decline unless full beneficial owner information and risk justifications are disclosed and verified.
		\end{block}
	\end{frame}
	
	% ================================================================
	% SECTION 32: ESG & AML INTEGRATION
	% ================================================================
	\begin{frame}
		\frametitle{\fontsize{13}{15}\selectfont \textbf{Section 32: ESG \& AML Integration}}
		\vspace{0.3cm}
		\begin{center}
			{\fontsize{12}{14}\selectfont \textbf{ESG-Aligned Compliance Requirements}}
		\end{center}
		\vspace{0.1cm}
		\begin{block}{\textbf{ESG Reporting Examples}}
			\scriptsize
			\begin{itemize}
				\item Publishing carbon emissions by facility
				\item Reporting executive salary data across genders
				\item Reporting supply chain labor standards
			\end{itemize}
		\end{block}
		\begin{block}{\textbf{UK-Based Institution Processing Payments}}
			\scriptsize
			Must conduct enhanced due diligence and evaluate the AML framework of the EEA-based partner before proceeding.
		\end{block}
	\end{frame}
	
	% ================================================================
	% EXAM TIPS SLIDE
	% ================================================================
%	\begin{frame}
%		\frametitle{\fontsize{13}{15}\selectfont \textbf{Exam Tips — Common Traps on CAMS Domain A}}
%		\vspace{0.1cm}
%		\begin{block}{\fontsize{9}{11}\selectfont \textbf{The Most Frequently Missed Distinctions:}}
%			\vspace{0.05cm}
%			\scriptsize
%			\begin{itemize}
%				\item \highlight{TF vs ML}: TF funds can be \textbf{legitimate}; ML funds are \textbf{always illicit}. TF = destination/intent; ML = source/disguise.
%				\item \highlight{SAR vs STR}: US BSA SARs are for defined financial institutions. DNFBPs outside US file STRs with national FIU.
%				\item \highlight{FATF Rec 15 vs 16}: Rec 15 = VASP framework. Rec 16 = Travel Rule. Both apply, but Rec 15 is the foundation.
%				\item \highlight{Domestic PEPs}: Require risk-based assessment. EDD if elevated risk.
%				\item \highlight{50\% Rule}: Own 50\%+ of an entity = entity is blocked, even if not listed.
%				\item \highlight{Structuring}: Illegal even with clean money. Breaking up cash to avoid CTR = crime in itself.
%				\item \highlight{Tipping Off}: Absolute criminal prohibition. Cannot disclose SAR filing to subject.
%				\item \highlight{SDD}: Not zero CDD. Any suspicion overrides SDD.
%				\item \highlight{Shell Banks}: Absolute prohibition. No EDD can make it acceptable.
%				\item \highlight{De-Risking}: FATF opposes blanket termination. Risk must be managed.
%				\item \highlight{SOW vs SOF}: Source of Wealth $\neq$ Source of Funds. FATF Rec 12 requires BOTH.
%				\item \highlight{PF vs TF}: PF = WMD $\rightarrow$ Rec 7. TF = terrorist acts $\rightarrow$ Rec 5.
%			\end{itemize}
%		\end{block}
%		\vspace{0.03cm}
%		\tipbox{\scriptsize \textbf{Final Tip:} For every answer choice, ask: \textit{``Is this always true, sometimes true, or never true?''} CAMS distractors are typically statements that are \textit{partially} correct. The correct answer is the one that is \textit{precisely} correct for the specific scenario.}
%	\end{frame}
	
	% ================================================================
	% SUMMARY TABLE
	% ================================================================
	\begin{frame}
		\frametitle{\fontsize{12}{14}\selectfont \textbf{Summary: Key Rules to Memorize for CAMS}}
		\vspace{0.1cm}
		\scriptsize
		\begin{longtable}{|p{0.28\textwidth}|p{0.67\textwidth}|}
			\hline
			\keypoint{Rule} & \keypoint{What the CAMS Exam Tests} \\
			\hline
			SAR vs STR & US BSA: defined FIs file SARs with FinCEN. US lawyers/real estate agents = NO BSA SAR. Non-US DNFBPs file STRs with national FIU. \\
			\hline
			TF vs ML & TF funds can be \textbf{legitimate}; intent/destination is key. ML funds always illicit; source concealment is key. \\
			\hline
			Sanctions & Strict liability (civil). OFAC 50\% rule: 50\%+ owned by designated = blocked even if unlisted. \\
			\hline
			Facilitation Payments & UK Bribery Act = zero tolerance, global. FCPA = narrow exception. \\
			\hline
			PEPs & Foreign PEPs: mandatory EDD. Domestic PEPs: risk-based assessment required. SOW and SOF both required. \\
			\hline
			VASPs & Foundational framework = Rec 15. Travel Rule = Rec 16. Both originator AND beneficiary required $\geq$\$1K. \\
			\hline
			Structuring & Illegal even with clean money. Crime = avoiding CTR filing. File SAR; do NOT tip off. \\
			\hline
			Correspondent Banking & EDD + senior approval required. Shell banks: absolute prohibition. \\
			\hline
			Beneficial Ownership & Natural person who ultimately owns/controls. 25\% = starting point only. \\
			\hline
			FATF Evaluations & Two dimensions: technical compliance + effectiveness. Gray List = EDD (Rec 19). \\
			\hline
		\end{longtable}
	\end{frame}
	
	% ================================================================
	% END SLIDE
	% ================================================================
	\begin{frame}
		\centering
		\vspace{1.5cm}
		{\fontsize{18}{22}\selectfont \textbf{Domain A — Complete Study Guide}}\\[0.3cm]
		{\fontsize{12}{14}\selectfont All answers verified against FATF Recommendations, US BSA, UK Bribery Act, and CAMS Study Guide principles.}\\[0.3cm]
		\rule{4cm}{0.5pt}\\[0.2cm]
		{\fontsize{9}{11}\selectfont \textit{For each question, understand WHY the wrong answers are wrong.} \\
			\textit{Exam distractors test the same misconceptions repeatedly.}}
		\vspace{0.5cm}
	\end{frame}
	
	% ================================================================
	% END DOCUMENT
	% ================================================================
\end{document}